%% bare_jrnl.tex
%% V1.4b
%% 2015/08/26
%% by Michael Shell
%% see http://www.michaelshell.org/
%% for current contact information.
%% 
%% This is a skeleton file demonstrating the use of IEEEtran.cls
%% (requires IEEEtran.cls version 1.8b or later) with an IEEE
%% journal paper.
%% 
%% Support sites:
%% http://www.michaelshell.org/tex/ieeetran/
%% http://www.ctan.org/pkg/ieeetran
%% and
%% http://www.ieee.org/

%%*************************************************************************
%% Legal Notice:
%% This code is offered as-is without any warranty either expressed or
%% implied; without even the implied warranty of MERCHANTABILITY or
%% FITNESS FOR A PARTICULAR PURPOSE! 
%% User assumes all risk.
%% In no event shall the IEEE or any contributor to this code be liable for
%% any damages or losses, including, but not limited to, incidental,
%% consequential, or any other damages, resulting from the use or misuse
%% of any information contained here.
%% 
%% All comments are the opinions of their respective authors and are not
%% necessarily endorsed by the IEEE.
%% 
%% This work is distributed under the LaTeX Project Public License (LPPL)
%% ( http://www.latex-project.org/ ) version 1.3, and may be freely used,
%% distributed and modified. A copy of the LPPL, version 1.3, is included
%% in the base LaTeX documentation of all distributions of LaTeX released
%% 2003/12/01 or later.
%% Retain all contribution notices and credits.
%% ** Modified files should be clearly indicated as such, including  **
%% ** renaming them and changing author support contact information. **
%%*************************************************************************


% *** Authors should verify (and, if needed, correct) their LaTeX system  ***
% *** with the testflow diagnostic prior to trusting their LaTeX platform ***
% *** with production work. The IEEE's font choices and paper sizes can   ***
% *** trigger bugs that do not appear when using other class files.       ***                          ***
% The testflow support page is at:
% http://www.michaelshell.org/tex/testflow/



\documentclass[journal]{IEEEtran}
% 
% If IEEEtran.cls has not been installed into the LaTeX system files,
% manually specify the path to it like:
% \documentclass[journal]{../sty/IEEEtran}





% Some very useful LaTeX packages include:
% (uncomment the ones you want to load)


% *** MISC UTILITY PACKAGES ***
%
%\usepackage{ifpdf}
% Heiko Oberdiek's ifpdf.sty is very useful if you need conditional
% compilation based on whether the output is pdf or dvi.
% usage:
% \ifpdf
%   % pdf code
% \else
%   % dvi code
% \fi
% The latest version of ifpdf.sty can be obtained from:
% http://www.ctan.org/pkg/ifpdf
% Also, note that IEEEtran.cls V1.7 and later provides a builtin
% \ifCLASSINFOpdf conditional that works the same way.
% When switching from latex to pdflatex and vice-versa, the compiler may
% have to be run twice to clear warning/error messages.






% *** CITATION PACKAGES ***
%
%\usepackage{cite}
% cite.sty was written by Donald Arseneau
% V1.6 and later of IEEEtran pre-defines the format of the cite.sty package
% \cite{} output to follow that of the IEEE. Loading the cite package will
% result in citation numbers being automatically sorted and properly
% "compressed/ranged". e.g., [1], [9], [2], [7], [5], [6] without using
% cite.sty will become [1], [2], [5]--[7], [9] using cite.sty. cite.sty's
% \cite will automatically add leading space, if needed. Use cite.sty's
% noadjust option (cite.sty V3.8 and later) if you want to turn this off
% such as if a citation ever needs to be enclosed in parenthesis.
% cite.sty is already installed on most LaTeX systems. Be sure and use
% version 5.0 (2009-03-20) and later if using hyperref.sty.
% The latest version can be obtained at:
% http://www.ctan.org/pkg/cite
% The documentation is contained in the cite.sty file itself.







% *** GRAPHICS RELATED PACKAGES ***
%
%\ifCLASSINFOpdf
%  % \usepackage[pdftex]{graphicx}
%  % declare the path(s) where your graphic files are
%  % \graphicspath{{../pdf/}{../jpeg/}}
%  % and their extensions so you won't have to specify these with
%  % every instance of \includegraphics
%  % \DeclareGraphicsExtensions{.pdf,.jpeg,.png}
%\else
%  % or other class option (dvipsone, dvipdf, if not using dvips). graphicx
%  % will default to the driver specified in the system graphics.cfg if no
%  % driver is specified.
%  % \usepackage[dvips]{graphicx}
%  % declare the path(s) where your graphic files are
%  % \graphicspath{{../eps/}}
%  % and their extensions so you won't have to specify these with
%  % every instance of \includegraphics
%  % \DeclareGraphicsExtensions{.eps}
%\fi
% graphicx was written by David Carlisle and Sebastian Rahtz. It is
% required if you want graphics, photos, etc. graphicx.sty is already
% installed on most LaTeX systems. The latest version and documentation
% can be obtained at:
% http://www.ctan.org/pkg/graphicx
% Another good source of documentation is "Using Imported Graphics in
% LaTeX2e" by Keith Reckdahl which can be found at:
% http://www.ctan.org/pkg/epslatex
%
% latex, and pdflatex in dvi mode, support graphics in encapsulated
% postscript (.eps) format. pdflatex in pdf mode supports graphics
% in .pdf, .jpeg, .png and .mps (metapost) formats. Users should ensure
% that all non-photo figures use a vector format (.eps, .pdf, .mps) and
% not a bitmapped formats (.jpeg, .png). The IEEE frowns on bitmapped formats
% which can result in "jaggedy"/blurry rendering of lines and letters as
% well as large increases in file sizes.
%
% You can find documentation about the pdfTeX application at:
% http://www.tug.org/applications/pdftex





% *** MATH PACKAGES ***
%
%\usepackage{amsmath}
% A popular package from the American Mathematical Society that provides
% many useful and powerful commands for dealing with mathematics.
%
% Note that the amsmath package sets \interdisplaylinepenalty to 10000
% thus preventing page breaks from occurring within multiline equations. Use:
%\interdisplaylinepenalty=2500
% after loading amsmath to restore such page breaks as IEEEtran.cls normally
% does. amsmath.sty is already installed on most LaTeX systems. The latest
% version and documentation can be obtained at:
% http://www.ctan.org/pkg/amsmath






% *** SPECIALIZED LIST PACKAGES ***
%
%\usepackage{algorithmic}
% algorithmic.sty was written by Peter Williams and Rogerio Brito.
% This package provides an algorithmic environment fo describing algorithms.
% You can use the algorithmic environment in-text or within a figure
% environment to provide for a floating algorithm. Do NOT use the algorithm
% floating environment provided by algorithm.sty (by the same authors) or
% algorithm2e.sty (by Christophe Fiorio) as the IEEE does not use dedicated
% algorithm float types and packages that provide these will not provide
% correct IEEE style captions. The latest version and documentation of
% algorithmic.sty can be obtained at:
% http://www.ctan.org/pkg/algorithms
% Also of interest may be the (relatively newer and more customizable) 
% algorithmicx.sty package by Szasz Janos:
% http://www.ctan.org/pkg/algorithmicx





% *** ALIGNMENT PACKAGES ***
%
%\usepackage{array}
% Frank Mittelbach's and David Carlisle's array.sty patches and improves
% the standard LaTeX2e array and tabular environments to provide better
% appearance and additional user controls. As the default LaTeX2e table
% generation code is lacking to the point of almost being broken with
% respect to the quality of the end results, all users are strongly
% advised to use an enhanced (at the very least that provided by array.sty)
% set of table tools. array.sty is already installed on most systems. The
% latest version and documentation can be obtained at:
% http://www.ctan.org/pkg/array


% IEEEtran contains the IEEEeqnarray family of commands that can be used to
% generate multiline equations as well as matrices, tables, etc., of high
% quality.




% *** SUBFIGURE PACKAGES ***
%\ifCLASSOPTIONcompsoc
%  % \usepackage[caption=false,font=normalsize,labelfont=sf,textfont=sf]{subfig}
%\else
%  % \usepackage[caption=false,font=footnotesize]{subfig}
%\fi
% subfig.sty, written by Steven Douglas Cochran, is the modern replacement
% for subfigure.sty, the latter of which is no longer maintained and is
% incompatible with some LaTeX packages including fixltx2e. However,
% subfig.sty requires and automatically loads Axel Sommerfeldt's caption.sty
% which will override IEEEtran.cls' handling of captions and this will result
% in non-IEEE style figure/table captions. To prevent this problem, be sure
% and invoke subfig.sty's "caption=false" package option (available since
% subfig.sty version 1.3, 2005/06/28) as this is will preserve IEEEtran.cls
% handling of captions.
% Note that the Computer Society format requires a larger sans serif font
% than the serif footnote size font used in traditional IEEE formatting
% and thus the need to invoke different subfig.sty package options depending
% on whether compsoc mode has been enabled.
%
% The latest version and documentation of subfig.sty can be obtained at:
% http://www.ctan.org/pkg/subfig




% *** FLOAT PACKAGES ***
%
%\usepackage{fixltx2e}
% fixltx2e, the successor to the earlier fix2col.sty, was written by
% Frank Mittelbach and David Carlisle. This package corrects a few problems
% in the LaTeX2e kernel, the most notable of which is that in current
% LaTeX2e releases, the ordering of single and double column floats is not
% guaranteed to be preserved. Thus, an unpatched LaTeX2e can allow a
% single column figure to be placed prior to an earlier double column
% figure.
% Be aware that LaTeX2e kernels dated 2015 and later have fixltx2e.sty's
% corrections already built into the system in which case a warning will
% be issued if an attempt is made to load fixltx2e.sty as it is no longer
% needed.
% The latest version and documentation can be found at:
% http://www.ctan.org/pkg/fixltx2e


%
%\usepackage{stfloats}
% stfloats.sty was written by Sigitas Tolusis. This package gives LaTeX2e
% the ability to do double column floats at the bottom of the page as well
% as the top. (e.g., "\begin{figure*}[!b]" is not normally possible in
% LaTeX2e). It also provides a command:
%\fnbelowfloat
% to enable the placement of footnotes below bottom floats (the standard
% LaTeX2e kernel puts them above bottom floats). This is an invasive package
% which rewrites many portions of the LaTeX2e float routines. It may not work
% with other packages that modify the LaTeX2e float routines. The latest
% version and documentation can be obtained at:
% http://www.ctan.org/pkg/stfloats
% Do not use the stfloats baselinefloat ability as the IEEE does not allow
% \baselineskip to stretch. Authors submitting work to the IEEE should note
% that the IEEE rarely uses double column equations and that authors should try
% to avoid such use. Do not be tempted to use the cuted.sty or midfloat.sty
% packages (also by Sigitas Tolusis) as the IEEE does not format its papers in
% such ways.
% Do not attempt to use stfloats with fixltx2e as they are incompatible.
% Instead, use Morten Hogholm'a dblfloatfix which combines the features
% of both fixltx2e and stfloats:
%
% \usepackage{dblfloatfix}
% The latest version can be found at:
% http://www.ctan.org/pkg/dblfloatfix




%
%
%
%
%
%
%
%
%
%
%
%
%
%
%
%
%
%
%
%
%
%
%
%
%
%
%
%
%
%
%
%
%
%
%
%
%
%
%
%
%
%
%
%
%
%
%
%
%
%
%
%
%
%
%
%
%
%
%
%
%
%
%
%
%
%
%
%
%
%
%
%
%
%
%
%
%
%
%
%
%
%
%
%
%
%
%
%
%
%
%
%
%
%
%
%
%
%
%
%
%
%
%
%
%
%
%
%
%
%
%
%
%
%
%
%
%
%
%
%
%
%
%
%
%
%
%
%
%
%
%
%
%
%
%
%
%
%
%
%
%
%
%
%
%
%
%
%
%
%
%
%
%
%
%
%
%
%
%
%
%
%
%
%
%
%
%
%
%
%
%
%
%
%
%
%
%
%
%
%
%
%
%
%
%
%
%
%
%
%
%
%
%
%
%
%
%
%
%
%
%
%
%
%
%
%
%
%
%
%
%
%
%
%
%
%
%
%
%
%
%
%
%
%
%
%
%
%
%
%
%
%
%
%
%
%
%
%
%
%
%
%
%
%
%
%
%
%
%
%
%
%
%
%
%
%
%
%
%
%
%
%
%
%
%
%
%
%
%
%
%
%
%
%
%
%
%
%
%
%
%
%
%
%
%
%
%
%
%
%
%
%
%
%
%
%
%
%
%
%
%
%
%
%
%
%
%
%
%
%
%
%
%
%
%
%
%
%
%
%
%
%
%
%
%
%
%
%
%
%
%
%
%
%
%
%
%
%
%
%
%
%
%
%
%
%
%
%
%
%
%
%
%
%
%
%
%
%
%
%
%
%
%
%
%
%
%
%
%
%
%
%
%
%
%
%
%
%
%
%
%
%
%
%
%
%
%
%
%
%
%
%
%
%
%
%
%
%
%
%
%
%
%
%
%
%
%
%
%
%
%
%
%
%
%
%
%
%
%
%
%
%
%
%
%
%
%
%
%
%
%
%
%
%
%
%
%
%
%
%
%
%
%
%
%
%
%
%
%
%
%
%
%
%
%
%
%
%
%
%
%
%
%
%
%
%
%
%
%
%
%
%
%
%
%
%
%
%
%
%
%
%
%
%
%
%
%
%
%
%
%
%
%
%
%
%
%
%
%
%
%
%
%
%
%
%
%
%
%
%
%
%
%
%
%
%
%
%
%
%
%
%
%
%
%
%
%
%
%
%
%
%
%
%
%
%
%
%
%
%
%
%
%
%
%
%
%
%
%
%
%
%
%
%
%
%
%
%
%
%
%
%
%
%
%
%
%
%
%
%
%
%
%
%
%
%
%
%
%
%
%
%
%
%
%
%
%
%
%
%
%
%
%
%
%
%
%
%
%
%
%
%
%
%
%
%
%
%
%
%
%
%
%
%
%
%
%
%
%
%
%
%
%
%
%
%
%
%
%
%
%
%
%
%
%
%
%
%
%
%
%
%
%
%
%
%
%
%
%
%
%
%
%
%
%
%
%
%
%
%
%
%
%
%
%
%
%
%
%
%
%
%
%
%
%
%
%
%
%
%
%
%
%
%
%
%
%
%
%
%
%
%
%
%
%
%
%
%
%
%
%
%
%
%
%
%
%
%
%
%
%
%
%
%
%
%
%
%
%
%
%
%
%
%
%
%
%
%
%\ifCLASSOPTIONcaptionsoff
  % \usepackage[nomarkers]{endfloat}
% \let\MYoriglatexcaption\caption
% \renewcommand{\caption}[2][\relax]{\MYoriglatexcaption[#2]{#2}}
%\fi
% endfloat.sty was written by James Darrell McCauley, Jeff Goldberg and 
% Axel Sommerfeldt. This package may be useful when used in conjunction with 
% IEEEtran.cls'  captionsoff option. Some IEEE journals/societies require that
% submissions have lists of figures/tables at the end of the paper and that
% figures/tables without any captions are placed on a page by themselves at
% the end of the document. If needed, the draftcls IEEEtran class option or
% \CLASSINPUTbaselinestretch interface can be used to increase the line
% spacing as well. Be sure and use the nomarkers option of endfloat to
% prevent endfloat from "marking" where the figures would have been placed
% in the text. The two hack lines of code above are a slight modification of
% that suggested by in the endfloat docs (section 8.4.1) to ensure that
% the full captions always appear in the list of figures/tables - even if
% the user used the short optional argument of \caption[]{}.
% IEEE papers do not typically make use of \caption[]'s optional argument,
% so this should not be an issue. A similar trick can be used to disable
% captions of packages such as subfig.sty that lack options to turn off
% the subcaptions:
% For subfig.sty:
% \let\MYorigsubfloat\subfloat
% \renewcommand{\subfloat}[2][\relax]{\MYorigsubfloat[]{#2}}
% However, the above trick will not work if both optional arguments of
% the \subfloat command are used. Furthermore, there needs to be a
% description of each subfigure *somewhere* and endfloat does not add
% subfigure captions to its list of figures. Thus, the best approach is to
% avoid the use of subfigure captions (many IEEE journals avoid them anyway)
% and instead reference/explain all the subfigures within the main caption.
% The latest version of endfloat.sty and its documentation can obtained at:
% http://www.ctan.org/pkg/endfloat
%
% The IEEEtran \ifCLASSOPTIONcaptionsoff conditional can also be used
% later in the document, say, to conditionally put the References on a 
% page by themselves.




% *** PDF, URL AND HYPERLINK PACKAGES ***
%
%\usepackage{url}
% url.sty was written by Donald Arseneau. It provides better support for
% handling and breaking URLs. url.sty is already installed on most LaTeX
% systems. The latest version and documentation can be obtained at:
% http://www.ctan.org/pkg/url
% Basically, \url{my_url_here}.




% *** Do not adjust lengths that control margins, column widths, etc. ***
% *** Do not use packages that alter fonts (such as pslatex).         ***
% There should be no need to do such things with IEEEtran.cls V1.6 and later.
% (Unless specifically asked to do so by the journal or conference you plan
% to submit to, of course. )


% *** ADDITIONAL PACKAGES ***
\usepackage[utf8]{inputenc}
\usepackage[T1]{fontenc}
\usepackage{amsmath}
\usepackage{amsfonts}
\usepackage{amssymb}
\usepackage{graphicx}
\usepackage{cite}
\usepackage{algorithm}
\usepackage{algorithmic}
\usepackage{array}
\usepackage{multirow}
\usepackage{booktabs}
\usepackage{tabularx}
\usepackage{subfig}
\usepackage{tikz}
\usetikzlibrary{shapes.geometric,positioning,arrows.meta,calc}
\usepackage{xcolor}
\usepackage{pifont}

% *** HYPERREF PACKAGE - MUST BE LOADED LAST ***
% IEEE standard for electronic submission and online publication
\usepackage[hidelinks]{hyperref}  % hidelinks: removes colored boxes around links (suitable for print/online)
\usepackage{url}  % Better URL formatting

% PDF metadata configuration for IEEE submission
\hypersetup{
    bookmarksnumbered=true,
    pdftitle={GAN-Based Document Restoration with Frozen HTR Recognizer and Multi-Component Loss Optimization for Historical Manuscripts},
    pdfauthor={Jatniko Nur Mutaqin and I Gusti Bagus Baskara Nugraha},
    pdfsubject={Document Image Enhancement, Handwritten Text Recognition, Deep Learning},
    pdfkeywords={Generative Adversarial Networks, Document Restoration, HTR, Dual-Modal Discriminator, Loss Function Optimization},
    pdfproducer={LaTeX with IEEEtran},
    pdfcreator={pdfLaTeX},
    pdfstartview=FitH
}

% correct bad hyphenation here
\hyphenation{op-tical net-works semi-conduc-tor hand-writ-ten doc-u-ment re-sto-ra-tion con-vol-u-tion-a-l re-cur-rent neu-ral net-works recog-ni-tion de-grad-a-tion AR-chi-val ANRI UNESCO}


\begin{document}
% 
% paper title
\title{GAN-Based Document Restoration with Frozen HTR Recognizer and Multi-Component Loss Optimization for Historical Manuscripts}
% 
% author names and affiliations
\author{Jatniko~Nur~Mutaqin
        and~I~Gusti~Bagus~Baskara~Nugraha,~\IEEEmembership{Member,~IEEE}% <-this % stops a space
\thanks{Jatniko Nur Mutaqin is a graduate student in the Smart System Study Program, School of Electrical Engineering and Informatics, Bandung Institute of Technology, Bandung 40132, Indonesia, and works at the National Archives of the Republic of Indonesia, Jakarta, Indonesia (e-mail: 23523314@mahasiswa.itb.ac.id).}% <-this % stops a space
\thanks{I Gusti Bagus Baskara Nugraha is a lecturer at the School of Electrical Engineering and Informatics, Bandung Institute of Technology, Bandung 40132, Indonesia (e-mail: baskara@stei.itb.ac.id).}% <-this % stops a space
\thanks{Manuscript received [Date]; revised [Date].}}

% The paper headers
\markboth{IEEE Transactions on [Journal Name],~Vol.~XX, No.~X, [Month]~20XX}% 
{Mutaqin \MakeLowercase{\textit{et al.}}: GAN-Based Document Restoration with Frozen HTR Recognizer for Historical Manuscripts}

% Make the title area
\maketitle

% Abstract
\begin{abstract}
Historical handwritten documents suffer from various degradation artifacts that significantly impair Handwritten Text Recognition (HTR) systems. While conventional document restoration methods focus on visual quality, they often fail to preserve text readability for HTR tasks. This paper proposes an HTR-oriented document restoration framework based on Generative Adversarial Networks (GAN) with two key contributions: (1) integration of a frozen pre-trained HTR recognizer that provides stable text-aware gradients without joint training instabilities, and (2) systematic multi-component loss function optimization through ablation studies. Unlike conventional approaches that rely solely on image-level supervision, our method explicitly optimizes text readability through gradients from the frozen recognizer. We conduct quantitative evaluation on a realistic semi-synthetic dataset constructed from authentic 16th-18th century paleographic documents from the National Archives of the Republic of Indonesia (ANRI) with controlled degradation augmentation to enable paired ground-truth evaluation. On 712 test images, our framework achieves a Character Error Rate (CER) of 34.9\%, reducing CER from 83.4\% (degraded condition) with 58.2\% relative improvement, while maintaining high visual quality (PSNR 30.74~dB, SSIM 0.987). This performance approaches the theoretical maximum on clean images (CER 34.1\%). Qualitative validation on 15 real ANRI manuscripts with authentic historical degradation demonstrates effective generalization to damage patterns that cannot be replicated synthetically. Ablation studies reveal that a 4-component loss configuration (pixel, adversarial, perceptual, CTC) achieves optimal balance, while dual-modal discriminator exploration shows marginal contribution ($\Delta$PSNR +0.28~dB, $p>0.05$), confirming that recognizer freezing strategy and loss optimization are more critical than discriminator architectural complexity.
\end{abstract}

% Keywords
\begin{IEEEkeywords}
Generative Adversarial Networks, Document Restoration, Handwritten Text Recognition, Loss Function Optimization, Historical Document Processing, Deep Learning.
\end{IEEEkeywords}







% For peer review papers, you can put extra information on the cover
% page as needed:
% \ifCLASSOPTIONpeerreview
% \begin{center} \bfseries EDICS Category: 3-BBND \end{center}
% \fi
% 
% For peerreview papers, this IEEEtran command inserts a page break and
% creates the second title. It will be ignored for other modes.
\IEEEpeerreviewmaketitle

% ============================================================ 
% I. INTRODUCTION
% ============================================================ 
\section{Introduction}
\IEEEPARstart{H}{istorical} handwritten documents constitute invaluable cultural heritage, containing critical information about past events, governance, trade, and social structures. The National Archives of the Republic of Indonesia (ANRI) preserves an extensive collection of colonial-era documents from the 16th to 18th centuries, including materials recognized by UNESCO as part of the Memory of the World Register.

These documents, written in Dutch paleographic script during the VOC (Vereenigde Oostindische Compagnie) and Dutch East Indies periods, provide unique insights into Indonesia's historical development. However, these manuscripts have suffered severe degradation including (Figure~\ref{fig:degradation_examples}): ink bleed-through, iron gall ink corrosion, paper aging, physical damage, and fading ink.

Such degradation significantly impairs text readability and poses substantial challenges for Handwritten Text Recognition (HTR) systems.

These degradation artifacts severely disrupt HTR system performance, which is essential for large-scale digitization, transcription, search, indexing, and information extraction from historical archives. Despite substantial advances in HTR technology using deep learning approaches, particularly Convolutional Recurrent Neural Networks (CRNN) and transformer-based architectures, these systems remain highly sensitive to image quality and degradation levels.

% Figure 1: Examples of degradation types in historical documents
% Subfigures showing: (a) background noise, (b) ink bleed-through, (c) paper aging/yellowing,
% (d) physical damage, (e) ink degradation, (f) combined degradation

\begin{figure}[!t]
\centering
\subfloat[Ink Bleed-Through]{\label{fig:degradation_examples_noise}%
\includegraphics[width=1.7in]{/home/lambda_one/tesis/GAN-HTR-ORI/docRestoration/Paper/main/fig_degradation_examples/ink_bleed_through.png}}\hfill
\subfloat[Iron Gall Ink Corrosion]{\label{fig:degradation_examples_inkCorrotion}%
\includegraphics[width=1.7in]{/home/lambda_one/tesis/GAN-HTR-ORI/docRestoration/Paper/main/fig_degradation_examples/ink_corrosion.png}}\hfill
\subfloat[Paper Aging (Yellowing, Stains)]{\label{fig:degradation_examples_aging}%
\includegraphics[width=1.7in]{/home/lambda_one/tesis/GAN-HTR-ORI/docRestoration/Paper/main/fig_degradation_examples/paper_aging.png}}\hfill
\subfloat[Physical Damage (Tears, Wrinkles)]{\label{fig:degradation_examples_damage}%
\includegraphics[width=1.7in]{/home/lambda_one/tesis/GAN-HTR-ORI/docRestoration/Paper/main/fig_degradation_examples/physical_damage.png}}\hfill
\subfloat[Ink Degradation or Fading]{\label{fig:degradation_examples_inkDeg}%
\includegraphics[width=1.7in]{/home/lambda_one/tesis/GAN-HTR-ORI/docRestoration/Paper/main/fig_degradation_examples/ink_degradation.png}}\hfill
\subfloat[Combined Degradation]{\label{fig:degradation_examples_combined}%
\includegraphics[width=1.7in]{/home/lambda_one/tesis/GAN-HTR-ORI/docRestoration/Paper/main/fig_degradation_examples/combined_degradation.png}}
\caption{Examples of degradation in 16th-18th century ANRI documents: (a) ink bleed-through, (b) iron gall ink corrosion, (c) paper aging, (d) physical damage, (e) fading ink, (f) combined degradation. Images demonstrate the varied degradation levels that challenge HTR systems.}
\label{fig:degradation_examples}
\end{figure}

\subsection{Limitations of yang ada Approaches}

Degraded historical documents present unique challenges for HTR systems, as noise, stains, blur, and background degradation can significantly reduce recognition accuracy. Automated document restoration aims to enhance image quality before HTR processing, yet yang ada approaches have fundamental limitations:

First, conventional document restoration methods (including classical thresholding such as Otsu~\cite{otsu1979threshold} and Sauvola~\cite{sauvola2000adaptive}, as well as deep learning approaches based on U-Net and GAN) optimize purely visual metrics (PSNR, SSIM) without considering the impact on HTR readability. Consequently, models may produce visually appealing images while damaging text structure (e.g., broken strokes, missing character connections) that is critical for HTR accuracy.

Second, HTR integration in document restoration has been proposed by Souibgui \textit{et al.}~\cite{souibgui2021enhance} through joint training of generator and recognizer. While innovative, this approach uses a single-modal discriminator that only evaluates visual quality, and joint training can cause optimization instabilities as the generator and recognizer compete for conflicting gradients.

Third, an evaluation gap exists—most benchmarks (DIBCO series) only provide visual ground truth without text transcriptions, making HTR impact validation infeasible without additional annotations. Datasets for 16th-18th century paleographic documents with authentic degradation are limited.

This research addresses these limitations with the following contributions. Core contributions: (1) Frozen HTR recognizer strategy providing stable text-aware gradients without joint training instabilities ($p<0.001$); (2) Multi-component loss function optimization through systematic ablation studies identifying an optimal 4-component configuration (pixel, adversarial, perceptual, CTC). Methodological contributions: (3) Evaluation on authentic 16th-18th century paleographic manuscripts from the National Archives of Indonesia; (4) Ablation studies measuring each component's contribution, including the finding that dual-modal discriminator provides marginal contribution ($\Delta$PSNR +0.28~dB, $p>0.05$), confirming that training protocol and generator capacity are more critical than discriminator architectural complexity.

\subsection{Research Contributions}

This research proposes an HTR-oriented document restoration framework that explicitly optimizes text readability while maintaining visual quality. Based on empirical validation through systematic ablation studies, the research contributions are organized as follows:

Core Contributions:
\begin{enumerate}
    \item Frozen HTR Recognizer Strategy: Using a frozen pre-trained HTR recognizer provides stable text-aware gradients without joint training instabilities. Ablation studies demonstrate significant impact on convergence stability ($p<0.001$) and preservation of recognition capabilities (see Section~V-B).
    
    \item Multi-Component Loss Function Optimization: Identification of an optimal 4-component configuration (adversarial, pixel reconstruction, perceptual, CTC) through systematic ablation studies. This configuration achieves Pareto balance between visual quality (PSNR 30.74~dB) and HTR readability (CER 34.9\%), outperforming the 5-component configuration (see Section~V-B).
\end{enumerate}

Methodological Contributions and Exploratory Findings:
\begin{enumerate}
    \setcounter{enumi}{2}
    \item Systematic Ablation Studies: Evaluation of each architectural and loss component's contribution, providing practical guidance for component selection. Finding: simple configurations with proper alignment achieve better performance than over-complex configurations.
    
    \item Evaluation on Authentic Historical Documents: Quantitative evaluation on realistic semi-synthetic dataset constructed from authentic 16th-18th century ANRI paleographic handwriting with controlled degradation augmentation ($n$=712), enabling precise measurement with paired ground truth. Qualitative validation on real ANRI manuscripts with authentic historical degradation ($n$=15) demonstrates effective generalization to damage patterns that cannot be replicated computationally.
    
    \item Findings on Dual-Modal Discriminator: Exploration of discriminator architecture with CNN (visual) and LSTM (textual) pathways reveals marginal contribution to final metrics ($\Delta$PSNR = +0.28~dB, $\Delta$CER = --0.01\%, $p>0.05$, statistically insignificant). This finding confirms that the key factors are training protocol and frozen recognizer strategy, not discriminator architectural complexity (see analysis in Section~VI-A).
\end{enumerate}

\subsection{Research Hypothesis and Validation}

This research validates the hypothesis that a document restoration framework with frozen HTR recognizer integration and HTR-oriented loss function optimization produces statistically significant Character Error Rate (CER) reduction compared to degraded baseline conditions (H1). Systematic ablation studies measure each component's contribution, with key findings: (1) recognizer freezing strategy proves to be a critical factor for stability and convergence, (2) 4-component loss configuration achieves optimal balance, and (3) dual-modal discriminator shows marginal contribution ($p>0.05$). These findings confirm that system success primarily depends on recognizer freezing strategy and loss optimization, not discriminator architectural complexity. Complete results and statistical analysis are presented in Section~V.

\subsection{Paper Organization}

The paper is organized as follows: Section~II reviews related work in document enhancement, GANs for image restoration, and HTR systems. Section~III presents the proposed architecture including Enhanced U-Net generator, dual-modal discriminator (CNN+LSTM), frozen recognizer integration, and optimized loss function with adaptive balancing. Section~IV describes experimental setup, dataset, training methodology, and hyperparameter optimization strategies. Section~V presents comprehensive results including quantitative comparison with degraded baseline input and clean ground truth reference, systematic ablation studies to validate each component's contribution, and qualitative analysis on real historical documents. Section~VI discusses findings, theoretical and practical implications, limitations, and future directions. Finally, Section~VII concludes the paper with a summary of contributions and research impact.

% ============================================================ 
% II. RELATED WORK
% ============================================================ 
\section{Related Work}

This section reviews prior research in three interconnected areas: document image restoration and enhancement, generative adversarial networks for image restoration, and handwritten text recognition systems.

\subsection{Document Image Restoration and Enhancement}

Document restoration aims to recover degraded documents by removing noise, artifacts, and background degradation while preserving text readability for OCR/HTR systems or human reading. Unlike binarization which produces binary output (black-and-white), document restoration can generate grayscale or color output that retains pixel intensity information for better visual quality.

\subsubsection{Classical Methods: Binarization and Thresholding}

Early approaches to document enhancement focused on binarization using thresholding. Global binarization methods such as Otsu~\cite{otsu1979threshold} determine a single threshold value for the entire image by maximizing inter-class variance. While computationally efficient, global methods tend to be less effective on documents with varying illumination or complex degradation.

Local adaptive thresholding methods address this limitation by computing pixel-specific thresholds based on local neighborhoods. Sauvola and Pietikäinen~\cite{sauvola2000adaptive} proposed a widely-used method that combines local mean and standard deviation. Niblack~\cite{niblack1986introduction} and Wolf~\cite{wolf2004extraction} developed similar local approaches with different adaptation strategies. However, these thresholding-based methods produce binary output that loses pixel intensity information and may eliminate fine details such as thin strokes.

\subsubsection{Deep Learning Methods for Document Restoration}

Since the mid-2010s, with advances in deep learning and availability of large labeled datasets such as DIBCO~\cite{pratikakis2016hdibco}, document restoration approaches have shifted from binarization to grayscale enhancement that preserves intensity information. U-Net-based auto-encoders~\cite{ronneberger2015unet} and Generative Adversarial Networks (GANs) enable end-to-end learning of mappings from degraded to clean documents with potential to preserve fine details critical for HTR.

\subsection{Generative Adversarial Networks for Image Restoration}

GANs~\cite{goodfellow2016improved} have substantially transformed the field of image restoration with their ability to generate realistic outputs through adversarial training between generator and discriminator. Unlike auto-encoders that only minimize reconstruction loss, GANs can produce sharper details and more natural textures.

\subsubsection{Conditional GANs for Image-to-Image Translation}

Pix2pix~\cite{isola2017image} introduced the conditional GAN (cGAN) framework for image-to-image translation tasks, where the generator receives an input image and produces a corresponding output image. The PatchGAN discriminator in Pix2pix evaluates authenticity at local patch level, enabling better capture of high-frequency details. This framework has been successfully applied to various domains including colorization, style transfer, and semantic segmentation.

CycleGAN~\cite{zhu2017unpaired} extends this concept to unpaired image-to-image translation using cycle consistency loss, enabling training without paired data. While useful for domain adaptation, unpaired approaches are less suitable for document restoration where paired degraded-clean data is available and pixel-level accuracy is crucial for preserving text content.

\subsubsection{GANs for Document Restoration and Enhancement}

In the context of document restoration, several GAN approaches have been proposed:

DE-GAN~\cite{souibgui2020degan} uses a U-Net generator with PatchGAN discriminator for document enhancement tasks including deblurring, watermark removal, and background cleaning. While achieving high visual quality, DE-GAN focuses on optimizing pixel-level similarity and does not explicitly integrate HTR readability metrics in its loss function.

DocEnTr~\cite{souibgui2022docent} introduces a document enhancement transformer that combines self-attention mechanisms with adversarial training. This approach demonstrates capability to capture long-range dependencies, but has high computational cost and has not integrated explicit HTR guidance.

HTR-GAN (Souibgui et al.)~\cite{souibgui2021enhance} represents the most relevant work to this research, integrating an HTR recognizer within the GAN framework for document enhancement. The generator and recognizer are jointly trained with combined losses from adversarial, reconstruction, and CTC. The joint training approach potentially causes optimization instabilities as the recognizer and generator compete for gradients—a theoretical concern that motivates the frozen recognizer approach in this research.

Distinctions of Proposed Approach: This research fundamentally differs from HTR-GAN in: (1) Frozen recognizer—we use a frozen pre-trained recognizer for stable gradients, not joint training; (2) Dual-modal discriminator—our discriminator processes images (CNN) and text sequences (LSTM) in parallel with potential for visual-textual coherence evaluation (although ablation studies show marginal impact, see Section~VI); (3) Adaptive loss balancing—adaptive balancing mechanism that dynamically adjusts ratios between CTC loss and visual losses for optimal balance; (4) Enhanced architecture—residual blocks, attention gates, and spatial attention for better preservation of thin strokes.

\subsection{Handwritten Text Recognition Systems}

HTR systems aim to recognize handwritten text from document images. The evolution of HTR has progressed from classical HMM-based methods to modern deep learning architectures.

CRNN (Convolutional Recurrent Neural Network)~\cite{shi2016crnn} combines CNN for feature extraction with RNN (typically LSTM/GRU) for sequence modeling, trained with CTC loss for alignment-free recognition. This architecture has become the foundation for many modern HTR systems due to its ability to handle variable-length sequences without explicit segmentation.

Attention mechanisms have improved HTR performance by enabling models to focus on relevant regions when recognizing each character. Encoder-decoder architectures with attention~\cite{bluche2017scan} can learn character segmentation and alignment implicitly.

Transformer-based HTR~\cite{michael2019evaluating,diaz2021rethinking} leverages self-attention to capture long-range dependencies without recurrent connections. Vision transformers and hybrid CNN-Transformer architectures demonstrate state-of-the-art performance on various handwriting benchmarks, including historical and modern documents.

Relevance to This Research: HTR systems are sensitive to image quality—degradation such as noise, blur, or broken strokes can increase CER. This research uses a frozen pre-trained hybrid CNN-Transformer recognizer to extract recognition features that guide the generator to produce HTR-readable output. The frozen recognizer approach differs from joint training because: (1) stable gradients—recognizer weights do not change; (2) consistent evaluation—recognizer performance does not degrade during GAN training; (3) modular design—recognizer can be updated independently.

\subsection{Gap Analysis and Positioning}

Based on the literature review above, Table~\ref{table_comparison} summarizes this research's position relative to yang ada approaches in terms of visual quality, HTR awareness, discriminator architecture, and loss optimization. The identified gaps are: (1) exploration of dual-modal (visual + textual) architecture in the discriminator as architectural completeness for bilateral validation—although ablation studies show marginal impact on final metrics (see Section~VI), (2) absence of systematic loss weight optimization to balance visual quality and HTR readability, and (3) potential instabilities from joint recognizer-generator training that can be addressed with a frozen recognizer approach, as motivated by theoretical concerns discussed in Section~II. Empirical validation of each component's relative contribution is presented in comprehensive ablation studies (Section~V), with Gap \#2 (systematic loss weight optimization) showing the most significant impact on final performance.

\begin{table}[!t]
\renewcommand{\arraystretch}{1.3}
\caption{Comparison of Document Enhancement Approaches: evaluation of state-of-the-art methods based on visual quality, HTR awareness, dual-modal discriminator architecture, and loss function optimization. The proposed approach integrates all features for HTR-oriented document restoration, with relative contribution of each component validated through ablation studies (Sections V--VI).}
\label{table_comparison}
\centering
\begin{tabular}{|l|c|c|c|c|}
\hline
\textbf{Method} & \textbf{Visual} & \textbf{HTR} & \textbf{Dual} & \textbf{Loss} \\
 & \textbf{Quality} & \textbf{Aware} & \textbf{Disc.} & \textbf{Optim.} \\
\hline
Classical & Low & No & No & N/A \\
U-Net & High & No & No & No \\
DE-GAN & High & No & No & No \\
Souibgui et al. & High & Yes & No & No \\
\textbf{Proposed} & \textbf{High} & \textbf{Yes} & \textbf{Yes} & \textbf{Yes} \\
\hline
\end{tabular}
\end{table}

% ============================================================ 
% III. PROPOSED METHOD
% ============================================================ 
\section{Proposed Method}

This section presents the proposed HTR-oriented document restoration framework. We first provide an architectural overview, then detail each component: the U-Net generator, dual-modal discriminator, frozen HTR recognizer integration, and optimized multi-component loss function.

\subsection{Framework Overview}

To address the limitations of conventional restoration methods that ignore machine text readability, this research proposes an HTR-oriented GAN framework that explicitly integrates a text recognizer in the training loop. The proposed framework, illustrated in Figure~\ref{fig:architecture_overview}, operates on degraded handwritten document line images and produces clean, readable output optimized for both visual quality and HTR performance.

% HTR-oriented document restoration framework architecture diagram
\begin{figure*}[!t]
\centering
\includegraphics[width=0.95\textwidth,keepaspectratio]{/home/lambda_one/tesis/GAN-HTR-ORI/docRestoration/Paper/drawio/main_framework.png}

\caption{HTR-oriented document restoration framework with dual-modal discriminator and frozen HTR recognizer. Three main components: (1) Enhanced U-Net Generator for image restoration, (2) Dual-Modal Discriminator (CNN+BiLSTM) for bilateral visual-textual validation, (3) Frozen HTR Recognizer for HTR-aware gradients. Training uses curriculum learning with multi-component loss (pixel, perceptual, adversarial, recognition feature, HTR). Total trainable parameters: 108.3M (Generator 89.4M + Discriminator 18.9M), Frozen recognizer: 27.86M.}
\label{fig:architecture_overview}
\end{figure*}

The restoration process flow is as follows: A degraded image $I_{\text{deg}}$ is input to generator $G$ which produces a restored image $I_{\text{gen}}$. The dual-modal discriminator $D$ receives pairs $(I_{\text{gt}}, \text{text}_{\text{gt}})$ as real samples and $(I_{\text{gen}}, \text{text}_{\text{pred}})$ as fake samples for adversarial learning. In parallel, the frozen HTR recognizer $R$ extracts features from $I_{\text{gen}}$ to ensure text readability through CTC loss and recognition feature loss.

The framework consists of four main components:

\begin{enumerate}
    \item Generator (G): Enhanced U-Net architecture mapping degraded images $I_{\text{deg}}$ to restored images $I_{\text{gen}}$. The generator is equipped with Residual Dense Blocks and Attention Gates for preserving fine text details while removing background noise.
    
    \item Dual-Modal Discriminator (D): Discriminator with parallel CNN and BiLSTM pathways that assess visual quality and text coherence simultaneously. Bilateral cross-attention fusion mechanism enables visual-textual consistency validation to avoid artifacts that damage character structure.
    
    \item Frozen HTR Recognizer (R): Pre-trained Transformer-based recognizer with fixed weights providing HTR-aware gradients without catastrophic forgetting risk. The recognizer achieves CER 33.72\% on initial ANRI validation dataset ($n$=948), and CER 26.57\% (validation set, $n$=710) and 34.1\% (test set, $n$=712) on clean ground truth images in this study—differences reflect text complexity variation across dataset splits.
    
    \item Multi-Component Loss Function: Weighted combination of $\mathcal{L}_{\text{pixel}}$ ($\lambda=50$), $\mathcal{L}_{\text{adv}}$ ($\lambda=3$), $\mathcal{L}_{\text{perc}}$ ($\lambda=1$), $\mathcal{L}_{\text{rec\_feat}}$ ($\lambda=8$), and $\mathcal{L}_{\text{CTC}}$ ($\lambda=0.15$) optimized through curriculum learning to balance visual quality and text readability.
\end{enumerate}

Unlike HTR-GAN~\cite{souibgui2021enhance} which jointly trains the recognizer with the generator, our method uses a frozen recognizer for training stability. The dual-modal discriminator architecture serves as exploration for bilateral visual-textual validation not present in previous methods that rely only on single visual discriminators.

The training process adopts alternating adversarial strategy with three-phase curriculum learning (detailed in Section~III-F):

\begin{itemize}
    \item Discriminator update: $D$ is trained to distinguish real pairs $(I_{\text{gt}}, \text{text}_{\text{gt}})$ from generated pairs $(I_{\text{gen}}, \text{text}_{\text{pred}})$, with dual-modal supervision on visual and textual domains simultaneously.
    
    \item Generator update: $G$ is optimized to fool $D$ while minimizing $\mathcal{L}_{\text{pixel}}$ against ground truth, $\mathcal{L}_{\text{perc}}$ for perceptual structure preservation, $\mathcal{L}_{\text{rec\_feat}}$ for HTR feature consistency, and $\mathcal{L}_{\text{CTC}}$ for explicit text readability.
\end{itemize}

\subsection{Generator Architecture: Enhanced U-Net with Residual Blocks}

This research adopts an Enhanced U-Net architecture that integrates four techniques for high-quality document restoration. Unlike conventional U-Net that only uses standard convolutional blocks, the Enhanced architecture integrates: (1) Residual Dense Blocks (RDB)~\cite{zhang2018residual} for hierarchical feature extraction with better gradient flow, (2) Convolutional Block Attention Module (CBAM)~\cite{woo2018cbam} for adaptive feature recalibration on channel and spatial dimensions, (3) Multi-Scale Feature Pyramid (MSFP)~\cite{lin2017fpn} with dilated convolutions~\cite{chen2017deeplab} at the bottleneck for multi-resolution context aggregation, and (4) Attention Gates~\cite{oktay2018attention} in the decoder for selective focus on text-bearing regions while suppressing background noise amplification. This combination of techniques enables preservation of thin strokes and paleographic diacritics while removing complex degradations such as bleed-through, stains, and noise.

\subsubsection{Network Specification}

Table~\ref{tab:generator_arch} details the complete generator architecture. The encoder progressively downsamples the input image through five stages, extracting multi-scale features at resolutions from $1024 \times 128$ to $32 \times 4$ (format: width $\times$ height). Each encoder block consists of:

\begin{itemize}
    \item Residual Dense Block (RDB): Three densely-connected convolutional layers (kernel $3 \times 3$) with growth rate 32, producing rich features through local fusion from all preceding layers. Unlike standard ResNet~\cite{he2016deep} which has only one shortcut connection, RDB uses dense connections allowing each layer to access features from all previous layers, enhancing representation capacity and gradient flow.
    \item Convolutional Block Attention Module (CBAM): Sequential attention mechanism consisting of: (a) channel attention with reduction ratio 8, using shared MLP for global average pooling and global max pooling pathways, and (b) spatial attention with kernel $7 \times 7$, combining channel-wise average and max pooling to generate focused attention maps. CBAM enables the network to adaptively focus on informative feature channels and relevant spatial regions.
    \item Batch normalization (momentum 0.9, eps $10^{-5}$) and LeakyReLU activation ($\alpha=0.2$)
    \item Residual connections for effective gradient propagation in deep architectures
\end{itemize}

The bottleneck at resolution $4 \times 32$ (height $\times$ width) uses Multi-Scale Feature Pyramid (MSFP)~\cite{lin2017fpn} with three dilated convolution scales (dilation rates 1, 2, 4)~\cite{chen2017deeplab} to capture multi-resolution context, followed by 512-channel RDB and CBAM for global feature refinement. The decoder mirrors the encoder with five upsampling stages, each equipped with attention gates for text region focus:

\begin{equation}
\alpha = \sigma(W_g^T g + W_x^T x + b)
\end{equation}

\noindent where $g \in \mathbb{R}^{H \times W \times C}$ is the decoder feature from coarser scale (gating signal), $x \in \mathbb{R}^{H \times W \times C}$ is the skip connection feature from the encoder, $W_g$ and $W_x$ are learnable projection matrices, $b$ is bias, and $\sigma$ is sigmoid activation function. The attention coefficient $\alpha \in [0,1]^{H \times W \times C}$ is then element-wise multiplied with the skip connection $x$ to suppress irrelevant background features while enhancing text-bearing regions: $\hat{x} = \alpha \odot x$. This mechanism allows the decoder to adaptively select relevant high-resolution features from the encoder.

% Generator Architecture Table
\begin{table*}[t]
\centering
\caption{Generator Architecture: Enhanced U-Net with Residual Blocks and Attention Gates}
\label{tab:generator_arch}
\begin{tabular}{|l|l|c|c|r|}
\hline
\textbf{Stage} & \textbf{Layer} & \textbf{Output Shape} & \textbf{Kernel/Stride} & \textbf{Parameters} \\
\hline
Input & Input & $128 \times 1024 \times 1$ & - & - \\
\hline
\multirow{2}{*}{Encoder-0} 
& ResBlock-1 & $128 \times 1024 \times 64$ & $3 \times 3 / 1$ & 4.7K \\
& MaxPool-1 & $64 \times 512 \times 64$ & $2 \times 2 / 2$ & - \\
\hline
\multirow{2}{*}{Encoder-1}
& ResBlock-2 & $64 \times 512 \times 128$ & $3 \times 3 / 1$ & 147.6K \\
& MaxPool-2 & $32 \times 256 \times 128$ & $2 \times 2 / 2$ & - \\
\hline
\multirow{2}{*}{Encoder-2}
& ResBlock-3 & $32 \times 256 \times 256$ & $3 \times 3 / 1$ & 590.1K \\
& MaxPool-3 & $16 \times 128 \times 256$ & $2 \times 2 / 2$ & - \\
\hline
\multirow{2}{*}{Encoder-3}
& ResBlock-4 & $16 \times 128 \times 512$ & $3 \times 3 / 1$ & 2.4M \\
& MaxPool-4 & $8 \times 64 \times 512$ & $2 \times 2 / 2$ & - \\
\hline
\multirow{2}{*}{Encoder-4}
& ResBlock-5 & $8 \times 64 \times 512$ & $3 \times 3 / 1$ & 2.4M \\
& MaxPool-5 & $4 \times 32 \times 512$ & $2 \times 2 / 2$ & - \\
\hline
Bottleneck & MSFP + RDB + CBAM & $4 \times 32 \times 512$ & - & 5.1M \\
\hline
\multirow{3}{*}{Decoder-4}
& UpSample-5 & $8 \times 64 \times 512$ & - & - \\
& AttentionGate-5 & $8 \times 64 \times 512$ & - & 262K \\
& RDB + CBAM & $8 \times 64 \times 512$ & $3 \times 3 / 1$ & 3.5M \\
\hline
\multirow{3}{*}{Decoder-3}
& UpSample-4 & $16 \times 128 \times 512$ & - & - \\
& AttentionGate-4 & $16 \times 128 \times 512$ & - & 262K \\
& RDB + CBAM & $16 \times 128 \times 512$ & $3 \times 3 / 1$ & 3.5M \\
\hline
\multirow{3}{*}{Decoder-2}
& UpSample-3 & $32 \times 256 \times 256$ & - & - \\
& AttentionGate-3 & $32 \times 256 \times 256$ & - & 66K \\
& RDB + CBAM & $32 \times 256 \times 256$ & $3 \times 3 / 1$ & 885K \\
\hline
\multirow{3}{*}{Decoder-1}
& UpSample-2 & $64 \times 512 \times 128$ & - & - \\
& AttentionGate-2 & $64 \times 512 \times 128$ & - & 16.5K \\
& RDB + CBAM & $64 \times 512 \times 128$ & $3 \times 3 / 1$ & 221K \\
\hline
\multirow{3}{*}{Decoder-0}
& UpSample-1 & $128 \times 1024 \times 64$ & - & - \\
& AttentionGate-1 & $128 \times 1024 \times 64$ & - & 4.1K \\
& RDB + CBAM & $128 \times 1024 \times 64$ & $3 \times 3 / 1$ & 55K \\
\hline
Output & Conv2D + Tanh & $128 \times 1024 \times 1$ & $1 \times 1 / 1$ & 65 \\
\hline
\multicolumn{4}{|r|}{\textbf{Total Parameters:}} & \textbf{21.8M} \\
\hline
\end{tabular}
\end{table*}

% Enhanced U-Net Generator Architecture Diagram
\begin{figure}[!t]
\centering
\includegraphics[width=0.5\textwidth,keepaspectratio]{/home/lambda_one/tesis/GAN-HTR-ORI/docRestoration/Paper/drawio/arsitektur_generator.png}
\caption{Enhanced U-Net Generator Architecture: 5-level encoder (RDB+CBAM), MSFP bottleneck (512 channels), decoder with attention gates. Skip connections preserve spatial details.}
\label{fig:unet_architecture}
\end{figure}

Architecture Statistics:
\begin{itemize}
    \item Total parameters: 21.8M (all trainable)
    \item Memory footprint: $\sim$87 MB (float32)
    \item Depth: 5 encoder stages + bottleneck + 5 decoder stages
    \item Input/Output: $1024 \times 128 \times 1$ (W$\times$H$\times$C)
    \item Output activation: Tanh (range [-1, 1])
\end{itemize}

Figure~\ref{fig:unet_architecture} visualizes the symmetric Enhanced U-Net architecture, where each encoder level has a decoder counterpart that communicates through skip connections with attention gates. Unlike conventional U-Net~\cite{ronneberger2015unet}, RDB+CBAM integration at each level increases representation capacity by $\sim$40\% without gradient degradation, while the MSFP bottleneck captures multi-scale global context before gradual reconstruction in the decoder.

\subsubsection{Design Rationale}

Residual Dense Blocks provide dual advantages: (1) better gradient propagation through multi-path dense connections, enabling training of deep architectures without gradient degradation, and (2) hierarchical feature reuse where each layer can access features from all previous layers, enhancing representation capacity without substantially increasing parameters. CBAM provides adaptive selectivity on channel and spatial dimensions, allowing the network to automatically focus on informative feature channels ("what") and relevant spatial regions ("where"), effective for suppressing background noise while enhancing text stroke features. Multi-Scale Feature Pyramid~\cite{lin2017fpn} at the bottleneck captures multi-resolution context through dilated convolutions~\cite{chen2017deeplab} with different scales, important for handling degradations of varying sizes (point noise vs. large stains). Attention Gates~\cite{oktay2018attention} provide spatial selectivity on skip connections, allowing the decoder to focus on text-bearing regions while suppressing noise amplification—crucial for preserving thin strokes and paleographic diacritics. The combination of these techniques yields an architecture capable of preserving fine details while removing complex degradations.

\subsection{Enhanced Dual-Modal Discriminator}

The dual-modal discriminator architecture evaluates images from both spatial (visual) and sequential (text) perspectives simultaneously. This design provides richer adversarial signals by considering both visual quality and textual coherence.

\subsubsection{Architecture Specification}

Table~\ref{tab:discriminator_arch} details the complete discriminator architecture. The discriminator accepts images ($128 \times 1024 \times 1$) and text sequences (maximum length 128 characters), processing them through two parallel pathways as illustrated in Figure~\ref{fig:discriminator_architecture}.

% Discriminator Architecture Table
\begin{table}[!t]
\centering
\caption{Enhanced Dual-Modal Discriminator Architecture: detailed specification of three branches (Image CNN, Text BiLSTM, Cross-Modal Fusion). Image branch processes visual features through 4 DownBlocks with spatial attention; text branch processes character sequences via embedding + BiLSTM + self-attention. Cross-modal fusion produces real/fake classification. Total: 17.4M parameters.}
\label{tab:discriminator_arch}
\scriptsize
\begin{tabular}{|l|l|c|r|}
\hline
\textbf{Branch} & \textbf{Layer} & \textbf{Output} & \textbf{Params.} \\
\hline
\multirow{9}{*}{\rotatebox{90}{\tiny CNN}}
& Input & $128{\times}1024{\times}1$ & - \\
& Conv2D+BN & $128{\times}1024{\times}64$ & 0.6K \\
& DownBlk-1 & $64{\times}512{\times}64$ & 37K \\
& DownBlk-2 & $32{\times}256{\times}128$ & 148K \\
& DownBlk-3 & $16{\times}128{\times}256$ & 590K \\
& DownBlk-4 & $8{\times}64{\times}512$ & 2.4M \\
& SpatialAttn & $8{\times}64{\times}512$ & 4.7K \\
& ResBlk-512 & $8{\times}64{\times}512$ & 2.4M \\
& AvgPool & $512$ & - \\
\hline
\multirow{5}{*}{\rotatebox{90}{\tiny LSTM}}
& Input & $128$ (seq.) & - \\
& Embed & $128{\times}128$ & 12.8K \\
& BiLSTM & $128{\times}512$ & 1.3M \\
& SelfAttn & $128{\times}512$ & 262K \\
& AvgPool & $512$ & - \\
\hline
\multirow{6}{*}{\rotatebox{90}{\tiny Fusion}}
& Proj-CNN & $128$ & 65.7K \\
& Proj-LSTM & $128$ & 65.7K \\
& CrossAttn & $256$ & 33K \\
& Dense-1+Drop & $512$ & 131K \\
& Dense-2+Drop & $256$ & 131K \\
& Sigmoid & $1$ & 257 \\
\hline
\multicolumn{3}{|r|}{\textbf{Total:}} & \textbf{17.4M} \\
\hline
\end{tabular}
\end{table}

% Dual-Modal Discriminator Architecture Diagram
\begin{figure}[!t]
\centering
\includegraphics[width=0.5\textwidth,keepaspectratio]{/home/lambda_one/tesis/GAN-HTR-ORI/docRestoration/Paper/drawio/arsitektur-dualModal.png}
\caption{Enhanced Dual-Modal Discriminator Architecture with bilateral cross-modal attention: CNN branch (512-D) with spatial attention, BiLSTM branch (512-D) with self-attention. Bilateral fusion (128-D) enables holistic validation.}
\label{fig:discriminator_architecture}
\end{figure}

Figure~\ref{fig:discriminator_architecture} visualizes the dual-modal architecture that distinguishes this approach from conventional discriminators. Unlike PatchGAN~\cite{isola2017image} which evaluates only visual patches, this architecture combines spatial analysis (CNN) and sequential analysis (BiLSTM) for holistic validation. The key characteristic is bilateral cross-modal attention (Img$\leftrightarrow$Text) enabling bidirectional interaction: the image branch validates text coherence from the HTR recognizer, while the text branch verifies visual stroke quality. This bilateral mechanism provides richer gradient feedback to the generator with 52\% parameter reduction (17.4M vs 137M baseline), driving simultaneous improvements in visual quality and textual consistency.

Image Branch (CNN with Spatial Attention): This branch extracts spatial visual features through a ResNet-based architecture. Input image $I \in \mathbb{R}^{128 \times 1024 \times 1}$ is processed through an initial convolutional layer (64 channels) with batch normalization and LeakyReLU activation ($\alpha=0.2$). Four residual downsampling blocks progressively reduce spatial resolution while increasing channel depth:
\begin{equation}
\begin{split}
128 \times 1024 \times 64 &\rightarrow 64 \times 512 \times 64 \rightarrow 32 \times 256 \times 128 \\
&\rightarrow 16 \times 128 \times 256 \rightarrow 8 \times 64 \times 512
\end{split}
\end{equation}
A spatial attention mechanism with $3 \times 3$ kernel is applied to the 512-channel features to focus on text-bearing regions:

\begin{equation}
\alpha_{\text{spatial}} = \sigma(\text{Conv}_{3 \times 3}(F_{\text{img}}))
\end{equation}

\noindent where $\sigma$ is the sigmoid function and $F_{\text{img}}$ is the feature map with dimensions $8 \times 64 \times 512$. Following spatial attention, an additional residual block with 512 channels (ResBlock-512) performs further feature processing without altering spatial dimensions, then global average pooling aggregates the $8 \times 64 \times 512$ spatial features into a representation vector $F_{\text{img}}^{\text{global}} \in \mathbb{R}^{512}$.

Text Branch (BiLSTM with Self-Attention): This branch captures sequential text coherence. Input character sequence $S \in \mathbb{N}^{L}$ (length $L \leq 128$) is first projected to embedding space $E \in \mathbb{R}^{L \times 128}$. BiLSTM processes the sequence bidirectionally to capture contextual dependencies:

\begin{equation}
h_t = [\overrightarrow{h_t}; \overleftarrow{h_t}] = [\text{LSTM}_{\text{fwd}}(E_t); \text{LSTM}_{\text{bwd}}(E_t)]
\end{equation}

\noindent producing hidden states $H \in \mathbb{R}^{L \times 512}$. Multi-head scaled dot-product self-attention is applied:

\begin{equation}
\text{Attention}(Q, K, V) = \text{softmax}\left(\frac{QK^T}{\sqrt{d_k}}\right)V
\end{equation}

\noindent where $Q, K, V$ are linear projections of $H$. Global average pooling yields text vector $F_{\text{text}}^{\text{global}} \in \mathbb{R}^{512}$.

Cross-Modal Fusion: Both feature vectors $F_{\text{img}}^{\text{global}} \in \mathbb{R}^{512}$ and $F_{\text{text}}^{\text{global}} \in \mathbb{R}^{512}$ are processed through a bilateral cross-modal attention mechanism to enable bidirectional interaction between modalities. The bilateral approach is necessary because: (1) images need to validate text coherence predicted by the HTR recognizer, and (2) text needs to verify the visual stroke quality that generates those characters. Thus, both modalities mutually provide context for more robust anomaly detection. Each modality projects its features to a low-dimensional shared space $d=128$ (reduced from 256 for stability) through six separate projection matrices:

Image$\rightarrow$Text Attention: Image as query, text as key/value:
\begin{align}
Q_{\text{img}} &= W_{\text{img}}^Q F_{\text{img}}^{\text{global}} \nonumber\\
K_{\text{text}} &= W_{\text{text}}^K F_{\text{text}}^{\text{global}} \nonumber\\
V_{\text{text}} &= W_{\text{text}}^V F_{\text{text}}^{\text{global}} \label{eq:cross_img_proj}\\
F_{\text{img}}^{\text{att}} &= \text{softmax}\left(\frac{Q_{\text{img}} K_{\text{text}}^T}{\sqrt{d}}\right) V_{\text{text}} \label{eq:cross_img_att}
\end{align}

\noindent where $W_{\text{img}}^Q, W_{\text{text}}^K, W_{\text{text}}^V \in \mathbb{R}^{512 \times 128}$ are learned projection matrices.

Text$\rightarrow$Image Attention: Text as query, image as key/value:
\begin{align}
Q_{\text{text}} &= W_{\text{text}}^Q F_{\text{text}}^{\text{global}} \nonumber\\
K_{\text{img}} &= W_{\text{img}}^K F_{\text{img}}^{\text{global}} \nonumber\\
V_{\text{img}} &= W_{\text{img}}^V F_{\text{img}}^{\text{global}} \label{eq:cross_text_proj}\\
F_{\text{text}}^{\text{att}} &= \text{softmax}\left(\frac{Q_{\text{text}} K_{\text{img}}^T}{\sqrt{d}}\right) V_{\text{img}} \label{eq:cross_text_att}
\end{align}

\noindent where $W_{\text{text}}^Q, W_{\text{img}}^K, W_{\text{img}}^V \in \mathbb{R}^{512 \times 128}$ are symmetric projection matrices

Both bilaterally-attended features ($F_{\text{img}}^{\text{att}}, F_{\text{text}}^{\text{att}} \in \mathbb{R}^{128}$) are concatenated into a fused vector: $F_{\text{fused}} = [F_{\text{img}}^{\text{att}}; F_{\text{text}}^{\text{att}}] \in \mathbb{R}^{256}$. The classification head consists of two dense layers (512 $\rightarrow$ 256 neurons) with dropout and LeakyReLU activation, followed by an output layer with sigmoid activation producing validity score $D(I, S) \in [0, 1]$ to distinguish genuine image-text pairs from generator-produced fake pairs.

Hyperparameter Configuration: Spatial attention kernel uses $3 \times 3$ size (reduced from $7 \times 7$ to minimize white-dot artifacts on strokes); cross-modal projection dimension set to 128 (reduced from 256 to improve training stability); LSTM units per direction are 256, yielding total BiLSTM output of 512 dimensions. Batch normalization momentum is 0.9 (increased from 0.8 for smoother gradient flow); dropout rate is 0.1 on classification head (reduced from 0.3 to preserve thin stroke information) and 0.2 on text embedding; activation function is LeakyReLU ($\alpha=0.2$) for CNN and sigmoid for LSTM recurrent activation (ensuring compatibility with standard LSTM implementations without cuDNN).

\subsubsection{Design Rationale}

Dual-Modal Approach: Text documents possess dual characteristics: (1) visual appearance (stroke quality, noise level), (2) sequential structure (character flow, inter-word spacing). The CNN pathway captures spatial quality, while BiLSTM captures sequential coherence. Bilateral cross-modal fusion enables the discriminator to assess image-text pairs holistically with bidirectional interaction, providing more informative adversarial signals to the generator for simultaneous improvements in visual quality and textual consistency.

Parameter Efficiency: The proposed discriminator achieves significant parameter efficiency (17.4M, 52\% reduction from 137M baseline) through simplified architecture without sacrificing discriminative capability: (1) low-dimensional projection (128) for cross-modal fusion reduces computational complexity while maintaining representation capacity; (2) smaller spatial attention kernel (3$\times$3) focuses on local details without artifacts; (3) optimized residual blocks balance architectural depth with training stability. Higher batch normalization momentum (0.9) yields more stable statistics and smoother gradient flow; lower dropout rate (0.1) preserves thin stroke information crucial for degraded document discrimination.

Capabilities: Through the combination of CNN spatial analysis and BiLSTM sequence processing, the discriminator can detect:
\begin{itemize}
    \item Broken or interrupted strokes disrupting text flow (via BiLSTM sequence anomaly detection)
    \item Noise patterns, artifacts, and bleed-through (via CNN spatial features)
    \item Inconsistent stroke width (via spatial attention mechanism)
    \item Unnatural gaps or connections between characters (via cross-modal fusion analysis)
\end{itemize}

By combining spatial and sequential assessment, the discriminator provides richer feedback to the generator, driving production of images that are both visually clean and textually coherent.

\subsection{Frozen HTR Recognizer Integration}

To explicitly optimize HTR performance, this research integrates a pretrained text recognizer with weight freezing strategy, where the recognizer is frozen (\texttt{model.trainable=False}) and used as a feature extractor.

\subsubsection{Architecture and Training Configuration}

The hybrid CNN-Transformer architecture employed consists of:
\begin{itemize}
    \item CNN Backbone: 4 convolutional blocks with progressive feature extraction, BatchNormalization, and GELU activation
    \item Sequence Projection: Dense layer + LayerNormalization + Dropout
    \item Transformer Encoder: 6 layers, 8 attention heads, positional encoding
    \item Output: CTC decoder for sequence prediction
\end{itemize}

The recognizer was trained on 4,739 Dutch handwritten line images from 16th--18th century ANRI colonial-era documents with 80/20 split (training $n$=3,791, validation $n$=948), achieving CER performance of 33.72\% on validation set. The model was then frozen (\texttt{trainable=False}) and integrated into the GAN training pipeline.

While CER 33.72\% is moderate, this performance is comparable to state-of-the-art on historical paleographic datasets (READ 2016: 25--35\%~\cite{sanchez2016icfhr}, ICDAR 2017: 28--42\%~\cite{sanchez2017icdar}) and reflects the intrinsic complexity of 16th--18th century handwriting. The HTR component functions as a feature extractor for text-aware guidance, not a visual optimizer.

Ablation study (Section~V-B, Table~\ref{table_ablation_loss}) validates the HTR role: at 15-epoch training, configuration with HTR (Experiment~4) achieved CER 29.66\%, while baseline without HTR (Experiment~1) could not measure CER as recognizer was not integrated. Degraded input showed CER 83.4\%. Production model with extended training (50 epochs) and HTR integration achieved CER 27.11\% on validation set ($n$=710) and CER 34.9\% on test set ($n$=712) with PSNR 30.74~dB. Note: The 7.79 percentage point CER difference between validation and test is not an overfitting indicator, but reflects intrinsic text complexity differences across splits\textemdash frozen recognizer on clean ground truth images also shows similar difference (26.57\% validation vs 34.1\% test, $\Delta$=7.53\%).

The primary HTR role is ensuring text preservation and readability (CER improvement 48.5--56.3 absolute points from degraded input on test and validation sets), while PSNR improvement stems from combined HTR guidance, adversarial training, perceptual loss, and extended training. This approach demonstrates method robustness to realistic HTR components without requiring perfect recognizer.

\subsubsection{Benefits of Freezing Strategy}

The recognizer freezing approach provides implementation advantages:
\begin{enumerate}
    \item Training stability: No joint optimization or conflicting gradients between generator and recognizer
    \item Computational efficiency: No backpropagation through recognizer, reducing computational burden
    \item Evaluation consistency: Fixed weights provide consistent CER/WER metrics throughout training
    \item Labeled supervision: Uses ground truth text for CTC loss supervision
\end{enumerate}

\subsubsection{Feature Extraction and Loss Computation}

This research extracts intermediate recognition features from the frozen recognizer:

\begin{equation}
\mathcal{F}_{\text{rec}}(I) = R_{\text{encoder}}(I)
\end{equation}

where $R_{\text{encoder}}$ represents the frozen encoder portion of the recognizer. Specifically, we use output from the projection layer before the transformer encoder, yielding feature maps with dimensions (batch, 128, 512). This layer was selected because it contains sequence-level features before contextual encoding, providing stable gradients for the generator. These features capture text-relevant characteristics such as stroke patterns, character shapes, and spatial arrangements.

Recognition feature loss between generated and ground truth images is computed using Mean Squared Error (MSE):

\begin{equation}
\begin{split}
\mathcal{L}_{\text{rec-feat}} &= \text{MSE}(\mathcal{F}_{\text{rec}}(I_{\text{gen}}), \mathcal{F}_{\text{rec}}(I_{\text{gt}})) \\
&= \|\mathcal{F}_{\text{rec}}(I_{\text{gen}}) - \mathcal{F}_{\text{rec}}(I_{\text{gt}})\|_2^2
\end{split}
\end{equation}

MSE was chosen because it is more sensitive to large feature deviations, providing stronger gradients for feature alignment improvement. This drives the generator to produce images whose text features match those of clean ground truth, implicitly enhancing HTR readability.

\subsection{Multi-Component Loss Function}

The generator is trained with a combination of five loss components calibrated through empirical grid search and validated on actual implementation:

\begin{equation}
\begin{split}
\mathcal{L}_{\text{total}} = &\lambda_{\text{adv}}\mathcal{L}_{\text{adv}} + \lambda_{\text{pixel}}\mathcal{L}_{\text{pixel}} + \lambda_{\text{perc}}\mathcal{L}_{\text{perc}} \\
&+ \lambda_{\text{ctc}}\mathcal{L}_{\text{ctc}} + \lambda_{\text{rec-feat}}\mathcal{L}_{\text{rec-feat}}
\end{split}
\end{equation}

Loss weight configuration for each component is shown in Table~\ref{table_optimal_weights}.

\subsubsection{Adversarial Loss}

Standard GAN adversarial loss drives realistic image generation:

\begin{equation}
\mathcal{L}_{\text{adv}} = -\log D(I_{\text{deg}}, I_{\text{gen}})
\end{equation}

Discriminator is trained to maximize:

\begin{equation}
\begin{split}
\mathcal{L}_D = -[&\log D(I_{\text{deg}}, I_{\text{gt}}) \\
&+ \log(1 - D(I_{\text{deg}}, I_{\text{gen}}))]
\end{split}
\end{equation}

\subsubsection{Pixel-Level Reconstruction Loss}

L1 loss ensures pixel-by-pixel similarity with ground truth:

\begin{equation}
\mathcal{L}_{\text{pixel}} = \|I_{\text{gen}} - I_{\text{gt}}\|_1
\end{equation}

L1 norm is chosen because it produces sharper results compared to L2 norm and is more robust to outliers.

\subsubsection{Perceptual Loss}

Perceptual loss~\cite{johnson2016perceptual} compares high-level features from pretrained VGG network:

\begin{equation}
\mathcal{L}_{\text{perc}} = \sum_{l \in L} \|\phi_l(I_{\text{gen}}) - \phi_l(I_{\text{gt}})\|_2
\end{equation}

where $\phi_l$ denotes features from layer $l$ of VGG-19 pretrained on ImageNet. Features are extracted from multiple convolutional layers for multiscale perceptual comparison. This approach ensures preservation of visual characteristics from low-level details (edges, textures) to high-level semantic features (stroke patterns, character shapes).

\subsubsection{Recognition Feature Loss}

As described in Section III-D, this component drives HTR-friendly characteristics using Mean Squared Error:

\begin{equation}
\mathcal{L}_{\text{rec-feat}} = \|\mathcal{F}_{\text{rec}}(I_{\text{gen}}) - \mathcal{F}_{\text{rec}}(I_{\text{gt}})\|_2^2
\end{equation}

% Table: Optimal loss weights from empirical tuning
\begin{table}[!t]
\renewcommand{\arraystretch}{1.3}
\caption{Loss Weight Configuration: optimal weights from empirical grid search to balance visual quality with text readability on validation dataset}
\label{table_optimal_weights}
\centering
\begin{tabular}{|l|c|c|}
\hline
Loss Component & Weight ($\lambda$) & Justification \\
\hline
Adversarial ($\mathcal{L}_{\text{adv}}$) & 3.0 & Enhanced realism \\
Pixel/L1 ($\mathcal{L}_{\text{pixel}}$) & 50.0 & Balanced preservation \\
Perceptual ($\mathcal{L}_{\text{perc}}$) & 1.0 & Topology preservation \\
Recognition Feature ($\mathcal{L}_{\text{rec-feat}}$) & 8.0 & Strong text guidance \\
CTC Loss ($\mathcal{L}_{\text{ctc}}$) & 0.15 & HTR monitoring (clipped) \\
\hline
\end{tabular}
\end{table}

Loss weight selection (Table~\ref{table_optimal_weights}) was performed through empirical tuning based on multi-objective trade-off analysis between visual quality, text preservation, and HTR readability. Dataset characteristics show dominance of thin strokes requiring specific balance between pixel reconstruction and text-based guidance.

High pixel reconstruction weight provides strong reconstruction signal while maintaining flexibility for adversarial enhancement, while moderate recognition feature weight provides text-aware guidance without dominating visual losses. Adaptive balancing mechanism is used for dynamic weight adjustment during training, preventing single loss component dominance and ensuring stable convergence.

\subsection{Training Strategy}

\subsubsection{Curriculum Learning and Optimization Protocol}

Training employs a three-phase curriculum learning approach for optimal stability and gradual convergence:

\begin{enumerate}
    \item Visual Warmup Phase (epochs 1--10): Exclusive training on visual quality without text supervision, allowing generator to build basic reconstruction capability. This phase prevents early mode collapse that can occur if CTC loss is introduced too early.
    
    \item Gradual CTC Ramp-up Phase (epochs 11--30): Linear gradual increase of CTC weight from zero to full weight (Table~\ref{table_optimal_weights}). This smooth integration allows generator to adapt to HTR gradients without disrupting established adversarial training stability.
    
    \item Full Training Phase (epochs 31--50): Training with all loss components active at optimal weights, accompanied by adaptive balancing to maintain dynamic equilibrium between visual and textual losses.
\end{enumerate}

Alternating adversarial training procedure with dual-modal supervision is applied:
\begin{enumerate}
    \item Discriminator step: Update $D$ to maximize ability to distinguish real pairs $(I_{\text{gt}}, \text{text}_{\text{gt}})$ from fake pairs $(I_{\text{gen}}, \text{text}_{\text{pred}})$ with simultaneous supervision on visual and textual domains.
    \item Generator step: Update $G$ to minimize $\mathcal{L}_{\text{total}}$ (combination of five loss components), producing images that simultaneously deceive discriminator and optimize HTR readability.
\end{enumerate}
Updates are performed with balanced 1:1 ratio, ensuring adversarial training stability.

Training configuration is based on empirical tuning for GAN stability and optimal convergence:

\begin{itemize}
    \item Optimization: Generator uses Adam algorithm with momentum parameters adjusted for GAN stability, while discriminator uses SGD with momentum to prevent mode collapse and ensure stable convergence. Both models apply gradient clipping to prevent gradient explosion.
    
    \item Learning rate: Both models use same learning rate with cosine annealing scheduling for smooth convergence without sharp drops.
    
    \item Batch size and dimensions: Small batch size with wide image dimensions preserving fine stroke details on single-line paleographic text.
    
    \item Early stopping: Applied with patience mechanism considering curriculum phases, restoring best weights based on combined PSNR-CER metric with weighting prioritizing visual quality.
\end{itemize}

\subsubsection{Stability Analysis and Regularization}

Figure~\ref{fig:convergence_analysis} shows training convergence trajectory over 50 epochs, validating stability and effectiveness of the proposed training strategy.

\begin{figure*}[!t]
\centering
\includegraphics[width=\textwidth]{/home/lambda_one/tesis/GAN-HTR-ORI/docRestoration/Paper/data_dukung/fig_convergence_analysis_combined.pdf}
\caption{Training convergence trajectory over 50 epochs: (a) Generator loss shows stable decrease without mode collapse, (b) PSNR increases reaching established targets, (c) SSIM achieves high values indicating structure preservation. Best model (marked by dashed vertical line and red star) selected based on combined PSNR-SSIM criteria. Smooth convergence curves validate effectiveness of proposed curriculum learning strategy.}
\label{fig:convergence_analysis}
\end{figure*}

Convergence Observations:
\begin{itemize}
    \item Generator loss: Smooth monotonic decrease in early phase, stable in late phase without oscillation, indicating stable adversarial training without mode collapse
    \item PSNR trajectory: Rapid increase in Visual Warmup Phase, followed by stable increase in Gradual CTC Ramp-up Phase, plateau in late phase triggering early stopping consideration
    \item SSIM progression: Consistent upward trend reaching high values at best epoch, demonstrating preservation of structural similarity alongside pixel-level accuracy
    \item Best model selection: Chosen based on combined PSNR-SSIM metric optimization, not arbitrary termination
    \item Statistical reliability: Narrow confidence intervals validate robust performance on validation set
\end{itemize}

Regularization Techniques: To prevent overfitting and improve model generalization, the following techniques were applied:

\begin{itemize}
    \item Batch normalization in all convolutional layers for training stability and accelerated convergence
    \item One-sided asymmetric label smoothing: applied only to real labels to reduce discriminator overconfidence while maintaining clear learning signal for fake samples
    \item Residual connections in discriminator for better gradient flow and stable deep architecture training
    \item Gradient clipping on both models (global norm 1.0) to prevent gradient explosion, specifically from CTC loss clipped at maximum value 400.0 for numerical stability
    \item Dropout on classification head (0.1) and embedding layer (0.2) for additional regularization
\end{itemize}

\paragraph{Technical Implementation Details} Model was trained with momentum 0.9 for SGD discriminator optimizer, batch size 2 on image dimensions $1024 \times 128 \times 1$ pixels, and Pure FP32 numerical precision (without mixed precision) for CTC loss stability. Early stopping uses patience 25 epochs with min delta 0.05 on combined PSNR-CER metric.

\paragraph{Early Stopping Note} Best model was selected based on combined PSNR-CER metric at epoch achieving optimal performance. Training continued to completion to verify no significant improvements occurred. Early stopping was not triggered because small fluctuations in validation metrics remained within specified threshold, confirming model reached stable convergence.

% ============================================================ 
% V. PENGATURAN EKSPERIMENTAL
% ============================================================ 
\section{Experimental Setup}

This section describes the datasets, baseline methods, evaluation metrics, implementation details, and hyperparameter optimization procedures used in the experiments.

\subsection{Datasets}

The proposed method is evaluated on a realistic semi-synthetic dataset constructed from authentic historical documents and qualitative validation on real manuscripts with natural degradation.

\subsubsection{Realistic Semi-Synthetic Dataset for Quantitative Evaluation}

Methodological Rationale: Quantitative evaluation of document restoration requires degraded-clean image pairs with perfect pixel-level alignment to accurately measure PSNR, SSIM, and CER metrics. Real historical documents with natural degradation lack clean ground truth as they cannot be restored to original condition. Therefore, a realistic semi-synthetic dataset was constructed by combining three authentic components~\cite{souibgui2021enhance,zhao2019document}:

Component 1 - Authentic Paleographic Handwriting (REAL):

Authentic 16th--18th century ANRI archival documents that are relatively clean (minimal visual degradation) are used as handwriting source. Documents annotated in Transkribus PAGE XML format were binarized using pretrained DE-GAN model~\cite{souibgui2020degan}, and individual text lines extracted using baseline coordinates with vertical padding, normalized to $128 \times 1024$ pixel dimensions. The extraction result is text line images with authentic paleographic characteristics\textemdash stroke width variation, complex ligatures, individual historical writer styles\textemdash that cannot be replicated by synthetic font rendering or generative models.

Component 2 - Authentic Historical Degradation Patterns (REAL):
Separate ANRI documents experiencing severe natural degradation (paper aging, organic stains, foxing, faded ink, iron gall ink corrosion) are used as degradation background texture sources. Background patches sized $128 \times 1024$ pixels are extracted from non-text regions of these damaged documents, preserving authentic historical degradation patterns formed over centuries of archival storage.

Component 3 - Controlled Degradation Augmentation (SEMI-SYNTHETIC):
To increase variation and create degraded-clean pairs with perfect pixel alignment, a layered degradation pipeline with six computational stages is stochastically applied to the combination of authentic text and real degradation backgrounds:

1. Texture Background Overlay:
Clean text image is composited onto real document texture patch extracted from damaged ANRI collection. Text regions are dimmed to simulate ink absorption into aged paper:
\begin{equation}
I_{\text{degradation}} = I_{\text{bg}} \cdot (1 - M_{\text{text}}) + (\alpha \cdot I_{\text{bg}}) \cdot M_{\text{text}}
\end{equation}
where $M_{\text{text}}$ is normalized text mask and $\alpha$ is dimming factor.

2. Bleed-Through Effect:
Simulates ink penetration from document reverse side by horizontally flipping random text image, applying Gaussian blur, and subtracting from background.

3. Perlin Noise-Based Organic Stains:
Generates natural-looking stain patterns using 3D Perlin noise~\cite{perlin1985image} with varied parameters to create smooth organic stain patterns.

4. Foxing Spots - Age Discoloration:
Simulates age-related discoloration by adding small brownish spots with random intensity.

5. Physical Damage:
Simulates document tears and folds through random lines and darkened patches.

6. Random Blur:
Applies Gaussian blur (smooth degradation), median blur (noise removal), or motion blur (scanner artifacts) with equal probability.

These effects are applied sequentially with randomized combinations, producing diverse and realistic degraded documents while maintaining perfect pixel-level ground truth alignment, eliminating manual annotation errors inherent in real document datasets.

Realistic Semi-Synthetic Dataset Statistics: The dataset consists of 4,739 degraded-clean image pairs ($128 \times 1024$ pixels, grayscale, normalized [0,1]) constructed from authentic ANRI documents. Each pair combines: (1) real historical paleographic handwriting, (2) degradation patterns from authentic naturally damaged archives, and (3) controlled augmentation for variation. This dataset preserves authentic historical document characteristics while providing paired ground truth necessary for high-precision quantitative evaluation.

% Figure 5 - Semi-synthetic realistic degradation-GT pair examples
\begin{figure}[!t]
\centering
\subfloat[Ground Truth Image - Authentic ANRI Paleographic Handwriting]{\label{fig:synthetic_gt}%

\includegraphics[width=3.2in]{/home/lambda_one/tesis/GAN-HTR-ORI/docRestoration/Paper/data_dukung/gambar_degGT/gt1.png}}\hfill
\subfloat[Degraded Image - Combination of Real Historical Degradation + Augmentation]{\label{fig:synthetic_deg}%
\includegraphics[width=3.2in]{/home/lambda_one/tesis/GAN-HTR-ORI/docRestoration/Paper/data_dukung/gambar_degGT/deg1.png}}
\caption{Realistic semi-synthetic dataset pair: (a) Ground truth from 16th--18th century ANRI documents showing authentic paleographic handwriting with stroke variation, ligatures, and historical writing styles that cannot be replicated by synthetic fonts, (b) Degraded image combines real historical degradation patterns (paper aging, organic stains from authentic archives) with controlled augmentation (bleed-through, blur), preserving realism of centuries-old document degradation while providing perfect pixel alignment for quantitative evaluation.}
\label{fig:synthetic_examples}
\end{figure}

\paragraph{Dataset Split (Academic Protocol)}
Data splitting follows academic three-way split protocol for unbiased evaluation:
\begin{itemize}
    \item Training: 70\% (3,317 images) - for generator and discriminator learning
    \item Validation: 15\% (710 images) - for hyperparameter tuning and early stopping  
    \item Test: 15\% (712 images) - held-out set for final evaluation
\end{itemize}

Dataset is shuffled with deterministic seed (seed=42) before splitting for reproducibility.

Degradation Stratification:
To ensure generalization, degradation pipeline is applied with controlled parameter randomization, ensuring model learns general degradation patterns and improves generalization capability on unseen historical documents.

\paragraph{Computational Specifications} Experiments were conducted on NVIDIA RTX A4000 (16GB VRAM), AMD Threadripper PRO 3955WX 16-core CPU, 128GB RAM, with TensorFlow 2.16.1, CUDA 12.8, and cuDNN 8.9.7. Training required approximately 48 GPU-hours for 50 epochs. Inference on single line image ($128\times1024$) requires 50--60ms, with batch processing throughput reaching $\approx$17 images/second (batch=2, 30--35ms per image) using 6--8GB VRAM.

\subsubsection{Real Historical Documents Dataset (ANRI)}

For generalization evaluation on real historical degradation, a representative document subset from National Archives of the Republic of Indonesia (ANRI) is used as qualitative evaluation set:

Test Document Selection: From ANRI archival collection of VOC and Dutch East Indies administrative documents (16th--18th centuries, primarily 1650--1800 era), 15 full-page documents were manually selected for generalization evaluation. Selection was based on criteria:
\begin{itemize}
    \item Degradation variation representation: Covers degradation spectrum from mild to severe to test model robustness
    \item Degradation type diversity: Paper aging, organic stains, foxing, faded ink, bleed-through, physical damage
    \item Paleographic characteristic variation: Various handwriting styles, contrast levels, and ligature complexity
\end{itemize}

\paragraph{Test Dataset Characteristics (n=15)}
Differences from Semi-Synthetic Dataset: This real ANRI dataset fundamentally differs from the semi-synthetic dataset for quantitative evaluation. These documents experienced authentic historical degradation during centuries of archival storage, displaying complex damage patterns that cannot be fully replicated through computational augmentation:
\begin{itemize}
    \item Temporal degradation: Natural paper aging with heterogeneous yellowing and local brittleness
    \item Organic artifacts: Biological stains, foxing spots with natural distribution patterns, non-uniform ink degradation
    \item Structural damage: Physical tears, permanent folds, material loss at document edges
    \item Writing variation: Dutch paleographic cursive handwriting with complex ligatures, period abbreviations, and writing style inconsistencies
\end{itemize}

Real ANRI documents lack paired visual ground truth (clean images) because authentic historical documents cannot be restored to original condition\textemdash physical and chemical damage over centuries is irreversible. Therefore, evaluation on real dataset is qualitative through visual inspection and paleography expert validation, complementing quantitative evaluation on semi-synthetic dataset. This dual evaluation approach follows standard practice in historical document restoration research~\cite{souibgui2021enhance,souibgui2020degan}: realistic semi-synthetic dataset (constructed from authentic documents with controlled augmentation) for precision quantitative measurement, and real dataset with natural degradation for qualitative generalization validation on real-world conditions.

\subsection{Baseline Methods and Evaluation Strategy}

Evaluation of the proposed method is conducted through two complementary approaches: (1) quantitative comparison with classical methods that can be replicated under identical experimental conditions, and (2) systematic ablation study to validate contribution of each architectural component.

\subsubsection{Classical Methods for Quantitative Comparison}

Classical methods are selected as quantitative baselines because they can be re-implemented on identical dataset without requiring retraining:

\begin{enumerate}
    \item No Enhancement (Degraded Input Baseline): Direct HTR on degraded images to measure degradation impact on HTR performance
    \item Otsu Thresholding~\cite{otsu1979threshold}: Histogram-based global binarization for document cleaning
    \item Sauvola Adaptive Thresholding~\cite{sauvola2000adaptive}: Local adaptive binarization for documents with non-uniform illumination
\end{enumerate}

Classical methods use standard OpenCV implementation with parameters optimized through grid search on validation subset.

\subsubsection{State-of-the-Art Deep Learning Methods (Research Context)}

Deep learning-based methods discussed in literature review (Section II) serve as research context to position the proposed contribution:

\begin{itemize}
    \item U-Net~\cite{ronneberger2015unet}: Encoder-decoder architecture with skip connections for image restoration
    \item DE-GAN~\cite{souibgui2020degan}: Conditional GAN for document enhancement with visual quality focus
    \item HTR-GAN~\cite{souibgui2021enhance}: GAN with jointly trained recognizer for HTR-oriented restoration
    \item DocEnTr: Transformer-based architecture with self-attention mechanism for document enhancement
\end{itemize}

Direct quantitative comparison with state-of-the-art deep learning methods is not performed because: (1) these methods are trained on different datasets with non-identical degradation characteristics, making direct comparison unfair~\cite{tensmeyer2017document}, (2) official reproducible implementations are not available for most methods, and (3) differences in architecture, training protocols, and loss functions make isolating specific contributions very difficult. This approach follows common evaluation practice in document restoration research where comparisons are performed only under identically replicable experimental conditions.

\subsubsection{Ablation Study as Component Validation}

To systematically validate contribution of each architectural component, incremental ablation study is performed by adding loss components progressively (Section V-B):

\begin{enumerate}
    \item Experiment 01 (Pixel Only): Equivalent to standard U-Net for image restoration
    \item Experiment 02 (+Adversarial): Equivalent to standard GAN without HTR component
    \item Experiment 03 (+Perceptual): Enhancement with VGG-based perceptual loss
    \item Experiment 04 (+CTC): Optimal configuration with HTR integration
    \item Experiment 05 (+Recognizer Features): Complete configuration for component redundancy validation
\end{enumerate}

This ablation approach provides comprehensive understanding of relative contribution of each component and empirically validates architectural design decisions. Ablation results (Table~\ref{table_ablation_loss}) show that 4-component configuration (Experiment 04) achieves optimal performance with best trade-off between visual quality and HTR readability.

\subsection{Evaluation Metrics}

Visual quality and text recognition metrics are used for evaluation.

\subsubsection{Visual Quality Metrics}

\begin{itemize}
    \item Peak Signal-to-Noise Ratio (PSNR): Measures pixel-level similarity
    \begin{equation}
    \text{PSNR} = 10 \log_{10}\left(\frac{\text{MAX}^2}{\text{MSE}}\right)
    \end{equation}
    where MAX=1.0 for images normalized to range [0,1], and MSE is Mean Squared Error.

    \item Structural Similarity Index (SSIM): Assesses perceived quality by considering luminance, contrast, and structure~\cite{wang2004ssim}
    \begin{equation}
    \text{SSIM}(x,y) = \frac{(2\mu_x\mu_y + c_1)(2\sigma_{xy} + c_2)}{(\mu_x^2 + \mu_y^2 + c_1)(\sigma_x^2 + \sigma_y^2 + c_2)}
    \end{equation}
\end{itemize}

\subsubsection{Artifact Detection Metrics (Secondary Monitoring)}

To monitor generation quality during training and ensure generator does not produce visual artifacts that interfere with text readability, three additional metrics are used:

\begin{itemize}
    \item \textit{Noise Variance}: Variance of Laplacian operator to measure high-frequency noise
    
    \item \textit{Isolated White Ratio}: Ratio of isolated white pixels to total pixels, detecting salt noise artifacts
    
    \item \textit{Local Variance}: Variance of local deviations from smoothed image, measuring restoration result smoothness
\end{itemize}

These metrics are used for training monitoring and quality validation, but not reported as primary evaluation metrics as they are diagnostic in nature.

\subsubsection{HTR Performance Metrics}

A hybrid CNN-Transformer model adapted on historical documents is used to evaluate HTR performance. This model is frozen during evaluation to ensure measurement consistency:

\begin{itemize}
    \item \textit{Character Error Rate} (CER): Percentage of character-level errors, calculated using Levenshtein edit distance:
    \begin{equation}
    \text{CER} = \frac{\text{EditDistance}(\text{prediction}, \text{reference})}{\text{length}(\text{reference})}
    \end{equation}
    
    \item \textit{Word Error Rate} (WER): Percentage of word-level errors, calculated with edit distance at word level:
    \begin{equation}
    \text{WER} = \frac{\text{EditDistance}(\text{predicted words}, \text{reference words})}{\text{length}(\text{reference words})}
    \end{equation}
\end{itemize}

\paragraph{Frozen Recognizer CER Baseline Clarification} To avoid confusion, this paper uses three different frozen recognizer CER baseline values:\vspace{-0.5em}
\begin{itemize}[itemsep=0pt]
    \item CER\textsubscript{pretrain} = 33.72\%: Performance of hybrid CNN-Transformer recognizer (CNN encoder with 6-layer transformer and CTC decoder) when evaluated on separate ANRI validation dataset ($n$=948) for initial generalization validation. This value is used to verify that recognizer can handle Dutch paleographic handwriting.
    \item CER\textsubscript{val-GT} = 26.57\%: Frozen recognizer performance on clean ground truth images from this study's validation set ($n$=710). This value represents theoretical upper bound for validation set.
    \item CER\textsubscript{test-GT} = 34.1\%: Frozen recognizer performance on clean ground truth images from this study's test set ($n$=712). This value represents theoretical upper bound for test set and serves as primary reference for assessing restoration quality.
\end{itemize}
\vspace{0.5em}
Difference between CER\textsubscript{val-GT} (26.57\%) vs CER\textsubscript{test-GT} (34.1\%), namely $\Delta$=7.53 percentage points, reflects that test set contains samples with higher paleographic complexity (more complex ligatures, greater handwriting style variation). This difference pattern is consistent with restoration results: CER\textsubscript{val-restored} = 27.11\% vs CER\textsubscript{test-restored} = 34.9\% ($\Delta$=7.79\%), confirming that gap is not caused by overfitting but rather dataset characteristic differences.

\begin{table}[!t]
\centering
\caption{Summary of CER Metrics Across Various Datasets and Conditions}
\label{table:cer_summary}
\small
\begin{tabular}{|l|c|c|c|}
\hline
\textbf{Condition} & \textbf{Dataset} & \textbf{$n$} & \textbf{CER (\%)} \\
\hline
\multicolumn{4}{|c|}{\textit{\textbf{Baseline Recognizer} on Clean Images}} \\
\hline
CER\textsubscript{pretrain} & ANRI initial validation & 948 & 33.72 \\
CER\textsubscript{val-GT} & Study validation GT & 710 & 26.57 \\
CER\textsubscript{test-GT} & Study test GT & 712 & 34.1 \\
\hline
\multicolumn{4}{|c|}{\textit{\textbf{Degraded Input} (No Restoration)}} \\
\hline
CER\textsubscript{degraded} & Validation \& Test & 1422 & 83.4 \\
\hline
\multicolumn{4}{|c|}{\textit{\textbf{Restoration Results} of Proposed Method}} \\
\hline
CER\textsubscript{val-restored} & Validation & 710 & 27.11 \\
CER\textsubscript{test-restored} & Test & 712 & \textbf{34.9} \\
\hline
\multicolumn{4}{|c|}{\textit{\textbf{Relative Improvement}}} \\
\hline
Validation & 83.4\% $\rightarrow$ 27.11\% & - & 67.5\% \\
Test & 83.4\% $\rightarrow$ 34.9\% & - & \textbf{58.2\%} \\
\hline
\end{tabular}
\end{table}

Table~\ref{table:cer_summary} summarizes all CER values reported in this paper for transparency and reproducibility. Use of CER\textsubscript{test-GT} = 34.1\% as primary reference throughout paper (abstract, discussion, conclusion) is consistent with standard evaluation protocol that uses test set as final metric.

\paragraph{Statistical Reporting and Evaluation Protocol}

All metrics are reported as mean $\pm$ standard deviation with 95\% confidence intervals, calculated from complete validation set (n=710) or test set (n=712). Statistical significance for model comparisons is tested using independent sample t-test with $\alpha = 0.05$. All experiments use fixed seed (seed=42) for full reproducibility.

Final Evaluation Protocol: All results reported in this paper (tables, abstract, hypothesis testing, discussion) use consistent evaluation on best checkpoint (epoch 44, selected based on validation metrics with early stopping patience=25) evaluated on held-out test set (n=712, never seen during training). This approach follows standard academic protocol to prevent overfitting and ensure unbiased results.

% ============================================================ 
% V. RESULTS AND DISCUSSION
% ============================================================ 
\section{Results and Discussion}

This section presents comprehensive experimental results including quantitative comparisons, ablation studies, qualitative analysis, and discussion of findings.

\subsection{Quantitative Results on Realistic Semi-Synthetic Test Set}

Table~\ref{table_synthetic_results} presents quantitative comparison results of the proposed method with classical methods on realistic semi-synthetic test set constructed from authentic ANRI documents. Evaluation focuses on classical baseline methods (Otsu thresholding and adaptive Sauvola) because other deep learning-based methods (U-Net, standard GAN, DE-GAN, HTR-GAN) are trained on different datasets making direct comparison unfair. This approach follows best practices in computer vision research where comparisons are performed only with methods that can be replicated under identical experimental conditions.

% Main quantitative results on synthetic test set
\begin{table*}[!t]
\renewcommand{\arraystretch}{1.3}
\caption{Quantitative Results on Realistic Semi-Synthetic Test Set (n=712, constructed from authentic ANRI documents)}
\label{table_synthetic_results}
\centering
\begin{tabular}{|l|c|c|c|c|c|}
\hline
\textbf{Method} & \textbf{PSNR (dB) $\uparrow$} & \textbf{SSIM $\uparrow$} & \textbf{F-measure $\uparrow$} & \textbf{CER (\%)} & \textbf{WER (\%)} \\
\hline
No Enhancement & 7.91$\pm$1.01 & 0.340$\pm$0.077 & -\textsuperscript{a} & 83.4$\pm$18.5 $\downarrow$ & 98.7$\pm$7.3 $\downarrow$ \\
Otsu Thresholding & 4.75$\pm$2.00 & 0.261$\pm$0.099 & 0.748$\pm$0.085 & 85.7$\pm$56.7 $\downarrow$ & 112.9$\pm$69.5 $\downarrow$ \\
Adaptive Sauvola & 9.92$\pm$1.99 & 0.437$\pm$0.149 & 0.933$\pm$0.030 & 109.5$\pm$179.5 $\downarrow$ & 171.5$\pm$172.2 $\downarrow$ \\
\hline
\textbf{Proposed} & \textbf{30.74$\pm$5.09} & \textbf{0.987$\pm$0.014} & \textbf{0.921$\pm$0.067} & \textbf{34.9$\pm$21.8} $\downarrow$ & \textbf{82.4$\pm$26.1} $\downarrow$ \\
\hline
\textit{Upper Bound (Clean GT)} & - & - & - & \textit{34.1$\pm$22.7} & \textit{82.1$\pm$27.3} \\
\hline
\end{tabular}
\textit{Note:} F-measure not applicable to degraded images without binarization. Classical methods use standard OpenCV implementation with optimal parameters. Proposed method uses epoch 44 checkpoint. All methods evaluated on identical test set ($n$=712) using same frozen recognizer. Clean GT Upper Bound (CER\textsubscript{test-GT} = 34.1\%) shows theoretical maximum recognizer performance on non-degraded images---see CER baseline clarification in Section~IV-C for explanation of different baseline values across datasets.
\end{table*}

Key observations based on quantitative results:
\begin{enumerate}
    \item Comparison with classical methods: The proposed method achieves PSNR 30.74$\pm$5.09 dB and SSIM 0.987$\pm$0.014, outperforming the best classical method (Sauvola: PSNR 9.92 dB, SSIM 0.437) by substantial margin. This result indicates superiority of deep learning-based approach for severely degraded document restoration.
    
    \item Substantial improvement on HTR metrics: CER decreased from 83.4\% (no enhancement) to 34.9\% (proposed method), showing absolute reduction of 48.5 percentage points or relative reduction of 58.2\%. Classical methods actually worsen HTR performance (Otsu: 85.7\%, Sauvola: 109.5\%) because inappropriate binarization produces artifacts and loss of important stroke details. This validates importance of HTR-oriented restoration approach over conventional binarization.
    
    \item Approaching theoretical upper bound: Proposed method CER (34.9\%) approaches clean GT upper bound (34.1\%), indicating that restoration successfully recovers readability close to original clean condition with difference of only 0.8 percentage points. This proximity shows that main HTR performance barrier is no longer image degradation, but rather limitations of HTR recognizer capability itself.
    
    \item Substantial quantitative improvement on visual metrics: Compared to degraded condition, proposed method improves PSNR by +22.83 dB (from 7.91 to 30.74 dB) and SSIM by +0.647 points (from 0.340 to 0.987), validating effectiveness of combining frozen HTR recognizer with HTR-oriented multi-component loss function optimization for simultaneous restoration of visual quality and text readability.
    
    \item Classical method limitations on complex degradation: Conventional thresholding methods fail on complex degradation because bimodal intensity distribution assumption does not hold for historical documents with spatially non-uniform degradation (foxing, bleed-through, heterogeneous aging). This is reflected in CER worse than no enhancement condition (Otsu: 85.7\% vs 83.4\%; Sauvola: 109.5\% vs 83.4\%).
    
    \item High variance in classical methods: Large CER standard deviation in Sauvola (179.5) and Otsu (56.7) indicates performance instability across samples with different degradation levels. Conversely, proposed method shows lower variance (21.8) thanks to adaptive learning from data enabling better generalization across various degradation levels.
\end{enumerate}

\subsubsection{Visual Method Comparison}

Figure~\ref{fig:method_comparison} presents qualitative visual comparison between proposed method and classical baseline methods on 5 representative samples from test set.

\begin{figure*}[!t]
\centering
\includegraphics[width=\textwidth]{/home/lambda_one/tesis/GAN-HTR-ORI/docRestoration/Paper/visualizations/fig_method_comparison_4methods.pdf}
\caption{Visual comparison of restoration methods on 5 representative test set samples. Rows show: (1) degraded input image, (2) Otsu result, (3) Sauvola result, (4) proposed method result, (5) clean ground truth. Proposed method produces high visual quality (PSNR 30.74~dB, SSIM 0.987) and higher HTR accuracy (CER 34.9\%) compared to classical methods that produce binarization artifacts.}
\label{fig:method_comparison}
\end{figure*}

Visual analysis reveals consistent performance patterns:
\begin{itemize}
    \item Restoration quality: Proposed method produces clean background with optimal stroke preservation, while Otsu produces salt-and-pepper artifacts and Sauvola tends to produce overly thick or broken strokes. This difference explains why classical methods actually worsen HTR performance despite being designed for document cleaning.
    \item Visual-HTR correlation: Samples with high PSNR/SSIM correlate with low CER, validating that visual quality improvement directly contributes to HTR accuracy improvement. This correlation justifies the research approach of simultaneously optimizing visual and HTR metrics.
    \item Performance consistency: The proposed method maintains consistent performance across samples with different degradation levels (see PSNR variation: 24.25--33.84 dB), while classical methods show inconsistency (CER varies drastically). This stability is important for large-scale archival restoration applications with high degradation variation.
\end{itemize}

\subsubsection{Validation Metric Trajectory Analysis}

Figure~\ref{fig:validation_metrics} presents detailed trajectories of four key validation metrics over 50 epochs with 95\% confidence intervals, showing model convergence process and statistical reliability of training results.

\begin{figure*}[!t]
\centering
\includegraphics[width=\textwidth]{/home/lambda_one/tesis/GAN-HTR-ORI/docRestoration/Paper/data_dukung/fig_validation_metrics_trajectory.pdf}
\caption{Validation metric trajectories with 95\% confidence intervals over 50 epochs (n=710): (a) PSNR reaches 30.91±0.42~dB, (b) SSIM reaches 0.987±0.001, (c,d) CER and WER approach clean image baseline. Best model (epoch 44, red star) selected based on combined PSNR-CER criteria. Narrow confidence intervals indicate consistent performance. Final evaluation on test set: PSNR 30.74~dB, SSIM 0.987, CER 34.9\% (Table~\ref{table_synthetic_results}).}
\label{fig:validation_metrics}
\end{figure*}

Validation trajectory observations confirming training stability:
\begin{enumerate}
    \item Gradual convergence on visual metrics: PSNR target achieved at epoch 38 (30.0 dB) and SSIM target at epoch 22 (0.95), indicating that visual warmup training phase (epochs 1--10) and gradual CTC increase phase (epochs 11--30) successfully built basic reconstruction capability before full HTR optimization.
    
    \item Confidence interval stability: Narrow confidence intervals for PSNR (±0.42 dB at epoch 44) and SSIM (±0.0011) show low variance on validation set, confirming model learns consistent restoration patterns without excessive fluctuation.
    
    \item Model selection based on combined criteria: Best model selected at epoch 44 (not final epoch) based on combined PSNR-CER criteria optimization on validation set, following early stopping protocol with patience=25. This selection prevents overfitting in late training phases.
    
    \item Generalization from validation to test: Performance on test set (PSNR 30.74~dB, SSIM 0.987, CER 34.9\%) consistent with validation trends, validating that model does not overfit. CER difference between validation and test (27.11\% vs 34.9\%, $\Delta$=7.79\%) consistent with frozen recognizer difference on clean images (26.57\% vs 34.1\%, $\Delta$=7.53\%), confirming gap is caused by higher intrinsic text complexity of test set, not model performance degradation. CER improvement ratio consistent: validation 67.5\% (83.4\%$\rightarrow$27.11\%), test 58.2\% (83.4\%$\rightarrow$34.9\%).
\end{enumerate}

These validation trajectories confirm effectiveness of the proposed three-phase curriculum learning strategy (Section III-F), where visual metric convergence occurs earlier (epochs 22--38) than HTR metrics, indicating that building visual reconstruction capability becomes foundation for text readability optimization in full training phase.

\subsection{Qualitative Evaluation on Real Historical Documents (ANRI)}

To evaluate model generalization to real historical degradation unseen during training, qualitative analysis was performed on documents from the National Archives of the Republic of Indonesia (ANRI). Due to limited reference transcriptions for historical paleographic documents, evaluation focuses on visual quality assessment based on systematic observation of restoration characteristics.

\subsubsection{Dataset and Evaluation Methodology}

Test documents: Dataset consists of 15 representative documents from ANRI archival collection from 17th--18th centuries (VOC and Dutch East Indies era, 1650--1800) with authentic degradation characteristics:
\begin{itemize}
    \item Paper aging: yellow/brown discoloration with high intensity
    \item Contrast variation: contrast range from low to high, covering severe to moderate degradation
    \item Foxing and organic stains: brownish spots from fungi and cellulose oxidation
    \item Bleed-through: ink transparency from verso page with varying intensity
    \item Iron gall ink artifacts: historical ink corrosion with dark halos and local pH changes
\end{itemize}

Dataset technical characteristics:
\begin{itemize}
    \item Resolution: 300 DPI from archival material scanning
    \item Orientation: portrait
    \item Degradation distribution: 2 severe documents (13.3\%), 8 moderate (53.3\%), 5 mild-moderate (33.3\%)
\end{itemize}
% bagaimana perintah tambah baris baru di latex



Inference procedure:
\begin{enumerate}
    \item Full-page documents processed using model with best validation performance
    \item Portrait-oriented overlap strategy to preserve paleographic feature authenticity
    \item Output format: uncompressed TIFF for archival preservation standards
    \item Visual evaluation: systematic comparison between degraded input and restored output for 5 quality aspects (background cleaning, text clarity, artifact presence, stroke preservation, overall utility)
\end{enumerate}

\subsubsection{Qualitative Results and Systematic Observations}

Figure~\ref{fig:anri_qualitative} shows representative restoration examples on authentic ANRI documents. Visual analysis reveals performance patterns consistent with input degradation characteristics.

% [FIGURE: ANRI Qualitative Results - 6 examples arranged in 2 rows]
\begin{figure*}[!t]
\centering
\includegraphics[width=7in]{/home/lambda_one/tesis/GAN-HTR-ORI/docRestoration/Paper/data_dukung/evaluasi_anri/fig_anri_qualitative_results.pdf}
\caption{Qualitative restoration on authentic 17th--18th century ANRI documents (n=15): each pair shows degraded input (left) and restored output (right). \textbf{Row 1 (a--c):} high-contrast success cases with substantial background cleaning and text preservation. \textbf{Row 2 (d--f):} challenging low-moderate contrast cases using conservative strategy to avoid artifacts.}
\label{fig:anri_qualitative}
\end{figure*}

\paragraph{Observations by Degradation Category (n=15 documents)}

\textit{High-Contrast Documents (n=5, 33.3\%):}
\begin{itemize}
    \item Background cleaning: foxing, aging stains, and organic spots effectively removed
    \item Text preservation: stroke thickness and character morphology preserved without excessive enhancement or thinning artifacts
    \item Bleed-through removal: verso text transparency minimized with good recto text preservation
    \item Artifacts: minimal artifacts, output suitable for transcription workflows
\end{itemize}

\textit{Moderate-Contrast Documents (n=8, 53.3\%):}
\begin{itemize}
    \item Partial success: background noise substantially reduced, but residual artifacts visible in regions with heavy bleed-through
    \item Text enhancement: character readability improved, but enhancement more conservative than high-contrast cases
    \item Trade-off: model prioritizes false negatives (missed noise) over false positives (stroke hallucination) to preserve authenticity
    \item Utility: output better than traditional binarization methods for moderate degradation
\end{itemize}

\textit{Low-Contrast Documents (n=2, 13.3\%):}
\begin{itemize}
    \item Limited enhancement: model adopts conservative strategy to avoid generating false artifacts
    \item Faded ink challenge: low-contrast ink receives only minimal enhancement
    \item Authenticity preservation: no hallucination or excessive enhancement that could mislead paleographic analysis
\end{itemize}

\subsubsection{Performance Analysis Based on Input Contrast}

To validate the hypothesis that restoration effectiveness correlates with input contrast level, quantitative analysis was performed on ANRI dataset degradation distribution.

\begin{table}[!t]
\renewcommand{\arraystretch}{1.3}
\caption{Restoration Performance on ANRI Documents by Input Contrast Level (n=15)}
\label{table:anri_contrast_analysis}
\centering
\small
\begin{tabular}{|l|c|c|c|}
\hline
\textbf{Contrast Category} & \textbf{n} & \textbf{\% Total} & \textbf{Observed Effectiveness} \\
\hline
High ($\geq$60) & 5 & 33.3\% & Excellent$^a$ \\
Moderate (30--60) & 8 & 53.3\% & Good$^b$ \\
Low ($<$30) & 2 & 13.3\% & Limited$^c$ \\
\hline
\multicolumn{4}{l}{\footnotesize $^a$Excellent: background cleaning >80\%, minimal artifacts} \\
\multicolumn{4}{l}{\footnotesize $^b$Good: substantial improvement, minor residual noise} \\
\multicolumn{4}{l}{\footnotesize $^c$Limited: conservative enhancement, no hallucination} \\
\end{tabular}
\end{table}

Contrast analysis findings show three key patterns:
\begin{enumerate}
    \item Contrast threshold effect: Model performance degradation becomes visible at contrast $<$30, indicating training dataset distribution limitations that may not fully cover low-contrast scenarios. This result is consistent with realistic semi-synthetic dataset contrast distribution having median contrast 35--45.
    
    \item Conservative strategy in ambiguity: In low SNR cases, model prioritizes precision over recall to avoid generating false artifacts---an important approach for historical archival preservation where document integrity is more important than maximal visual enhancement.
    
    \item Success on typical degradation: Of 15 ANRI documents, 13 documents (86.7\%) with contrast $\geq$30 showed useful improvement for transcription workflows, while 10 documents (66.7\%) showed substantial enhancement with background cleaning >70\% and minimal artifacts. This difference reflects restoration quality gradation: "useful" category includes all documents with minimal to optimal improvement, while "significant" is limited to documents with substantial background cleaning.
\end{enumerate}

Identified generalization limitations:
\begin{enumerate}
    \item Extreme fading: Documents with contrast $<$20 (1/15, 6.7\%) received minimal enhancement due to different semi-synthetic training dataset contrast distribution. This is consistent with validation set observations where performance declines at input PSNR $<$8 dB (region outside training distribution).
    
    \item Bleed-through variance: Real-world bleed-through patterns with complex verso page layouts are more challenging than synthetic augmentation using simple horizontal flip, resulting in partial success in 4/15 cases (26.7\%).
    
    \item Iron gall ink artifacts: Historical chemical-specific degradation (halo effects, pH-induced color shifts) not fully replicated in synthetic training, causing occasional over-suppression artifacts in 2/15 documents (13.3\%).
\end{enumerate}

Practical implications for archival digitization: ANRI evaluation results confirm that the model can generalize to real historical documents with degradation characteristics overlapping with synthetic training distribution (contrast $\geq$30, moderate to severe aging). For extreme cases outside training domain, model adopts conservative strategy that preserves document authenticity, avoiding catastrophic failure modes such as hallucination or structural corruption. This approach aligns with archival preservation principles where false negatives (missed restoration opportunities) are more acceptable than false positives (introduced artifacts that mislead historical interpretation).

\subsubsection{Expert Validation (Paleographic Assessment)}

To validate restoration quality from professional paleographic perspective, this study involved two independent paleography experts with over 10 years experience in analyzing historical documents from VOC and Dutch East Indies periods (16th--18th centuries).

Validation methodology:
\begin{itemize}
    \item Expert panel: 
    \begin{itemize}
        \item[$-$] \textit{Rater 1:} Paleography expert, specializing in 16th--18th century ancient Dutch and Latin scripts, 16 years experience in VOC document conservation (affiliation withheld for evaluation blindness)
        \item[$-$] \textit{Rater 2:} Historical archival restoration expert, specializing in Dutch East Indies administrative documents, experienced in physical archival restoration and collection digitization (affiliation withheld)
    \end{itemize}
    \item Blind assessment protocol: Independent evaluation on 15 document pairs (degraded-restored) without technical information about model or restoration methodology. Both raters were unaware that restoration was performed using GAN or HTR-based approach. Each document evaluated using 4-point Likert scale (1=Poor, 2=Fair, 3=Good, 4=Excellent) for four quality criteria.
    \item Evaluation criteria:
    \begin{enumerate}
        \item \textit{Text Readability:} Ease of reading paleographic characters, especially 17th century minuscule letters and ligatures
        \item \textit{Paleographic Authenticity:} Preservation of period handwriting characteristics (ductus, pen pressure, stroke thickness variation)
        \item \textit{Artifact Minimization} (inverse scoring): Level of restoration artifact interference with paleographic interpretation (high score = minimal artifacts)
        \item \textit{Transcription Utility:} Output suitability for archival digitization workflows and systematic transcription
    \end{enumerate}
    \item Inter-rater reliability: Cohen's kappa calculated to measure agreement consistency between two independent raters, with threshold $\kappa > 0.60$ for substantial agreement (Landis \& Koch, 1977)
\end{itemize}

\begin{table}[!t]
\renewcommand{\arraystretch}{1.3}
\caption{Paleographic Expert Validation Results on ANRI Documents (n=15)}
\label{table:anri_expert_validation}
\centering
\small
\begin{tabular}{|l|c|c|c|}
\hline
\textbf{Criteria} & \textbf{Rater 1} & \textbf{Rater 2} & \textbf{Mean} \\
\hline
Text Readability & 3.47 & 3.33 & 3.40 \\
Paleographic Authenticity & 3.60 & 3.53 & 3.57 \\
Artifact Minimization (inv.) & 3.27 & 3.13 & 3.20 \\
Transcription Utility & 3.53 & 3.47 & 3.50 \\
\hline
\textbf{Overall Score} & \textbf{3.47} & \textbf{3.37} & \textbf{3.42} \\
\hline
\multicolumn{4}{l}{\footnotesize Cohen's kappa (inter-rater): $\kappa$ = 0.78 (substantial agreement)} \\
\multicolumn{4}{l}{\footnotesize Scale: 1=Poor, 2=Fair, 3=Good, 4=Excellent} \\
\multicolumn{4}{l}{\footnotesize 95\% CI for overall mean: [3.28, 3.56]} \\
\end{tabular}
\end{table}

Score distribution per contrast category:

Stratification analysis shows strong correlation between input contrast and paleographic scores:

\begin{itemize}
    \item High Contrast ($\geq$60, n=5): Mean score 3.75/4.0. Both raters gave high appreciation for preservation of microscopic details such as hairline strokes and pen pressure modulation crucial for comparative paleographic analysis.
    
    \item Moderate Contrast (30--60, n=8): Mean score 3.38/4.0. Positive evaluation with minor notes that occasional slight blur occurs in regions with heavy bleed-through, which can reduce clarity of differentiation between recto and verso text, but does not impede main transcription.
    
    \item Low Contrast ($<$30, n=2): Mean score 2.63/4.0. Raters appreciated model's conservative strategy avoiding false reconstruction. Very low contrast documents judged to maintain paleographic integrity despite minimal visual enhancement.
\end{itemize}

Expert qualitative feedback:

\textit{Rater 1 (paleography expert):} ``Restoration shows implicit understanding of 17th century Dutch script characteristics, particularly preservation of Gothic minuscule letter aspect ratios and characteristic ligatures such as 'st', 'ct', and 'ck'. For high-contrast documents, foxing cleaning very effective without sacrificing superficial details such as pen lift marks important for authenticity analysis. Critical note: in 2--3 moderate cases, slight blur occurs in regions with corroded iron gall ink, but this does not reduce utility for semantic transcription.''

\textit{Rater 2 (archival restoration expert):} ``From digital archival perspective, model output highly suitable for mass digitization with minimal manual correction. Background noise reduction in 86.7\% documents (13/15) significantly facilitates downstream OCR/HTR. Special appreciation for conservative approach on extreme faded documents---model does not perform creative reconstruction that could mislead historical interpretation. For ANRI workflows, recommended for batch preprocessing use with expert verification on cases with contrast $<$30.''

Inter-rater agreement: Cohen's kappa = 0.78 indicates substantial agreement (Landis \& Koch, 1977), validating paleographic evaluation consistency. Greatest disagreement occurred on "Artifact Minimization" criteria for moderate-contrast documents (30--60), where Rater 1 was more tolerant of residual noise than Rater 2 who prioritized cleanliness for OCR automation. However, both raters agreed on relative document ranking (Spearman's $\rho$ = 0.91, p$<$0.001).

Implications for archival digitization: Expert validation confirms that model meets paleographic quality standards for three key applications: (1) Digital archival preservation---paleographic feature authenticity preserved (score 3.57/4.0), enabling comparative script analysis and scribe identification; (2) Public accessibility---readability improvement (score 3.40/4.0) reduces barriers for non-specialist researchers; (3) Semi-automated workflows---transcription utility (score 3.50/4.0) supports integration with modern HTR pipelines with minimal expert supervision.

\subsection{Degradation Level Correlation Analysis}

To understand the relationship between document degradation level and restoration effectiveness, comprehensive correlation analysis was performed using validation set (n=710 samples). This analysis answers key research question: \textit{``Is restoration more effective for heavily degraded or lightly degraded documents?''}

\begin{figure*}[!t]
\centering
\includegraphics[width=\textwidth]{/home/lambda_one/tesis/GAN-HTR-ORI/docRestoration/Paper/data_dukung/fig_combined_correlation_analysis.pdf}
\caption{Degradation level correlation analysis on validation set (n=710): (a) strong negative correlation (Pearson r=--0.963, p<0.001) between SNR and CER improvement shows restoration more effective on heavily degraded documents, (b) PSNR vs post-restoration CER shows model brings all degradation levels close to clean image reference value (26.6\%, green line), (c) SSIM vs CER shows weak correlation indicating structural similarity not sole predictor for HTR accuracy, (d) before-after CER comparison shows majority of samples experience improvement (points below diagonal y=x, gray). Note: Data from validation set (n=710) with statistics PSNR 30.91±5.60 dB, CER 27.11±20.83\%, SSIM 0.987±0.014. Final evaluation on test set (n=712) reported in Table~\ref{table_synthetic_results}: PSNR 30.74 dB, CER 34.9\%, SSIM 0.987.}
\label{fig:degradation_correlation}
\end{figure*}

\subsubsection{Correlation Findings}

\textit{SNR vs CER Improvement (Panel a):}
\begin{itemize}
    \item Very strong negative correlation: Pearson's r = --0.963, R² = 0.927, p < 0.001
    \item Interpretation: documents with low SNR (heavy degradation) achieve greater CER improvement than documents with high SNR (light degradation)
    \item Practical implication: model shows higher effectiveness for severely degraded historical documents, with CER improvement inversely proportional to initial SNR level
    \item Empirical validation: strong correlation (R² > 0.9) confirms that initial degradation level is significant predictor for restoration benefit
\end{itemize}

\textit{PSNR vs Restored CER (Panel b):}
\begin{itemize}
    \item Weak correlation: Pearson's r approaches zero
    \item Interpretation: after restoration, PSNR (image quality) does not strongly correlate with CER (HTR accuracy), indicating that visual quality improvement does not automatically guarantee text recognition accuracy improvement
    \item Practical implication: model tends to bring various degradation levels to relatively uniform HTR performance approaching clean baseline value
    \item Density visualization: concentration around baseline CER (26.6\%) confirms post-restoration HTR performance normalization
\end{itemize}

\textit{SSIM vs CER (Panel c):}
\begin{itemize}
    \item Weak correlation observed between SSIM and CER, indicating SSIM alone insufficient to predict HTR accuracy
    \item Insight: structural similarity (SSIM) measures global visual quality, but HTR performance depends on character-level details (stroke sharpness, local contrast) not fully captured by structural metrics
    \item Implication: justification for multi-metric evaluation (PSNR + SSIM + CER) rather than relying on single visual quality metric to validate HTR-oriented restoration effectiveness
\end{itemize}

\textit{Before vs After CER (Panel d):}
\begin{itemize}
    \item Majority of samples lie below diagonal line (y=x), confirming that restoration generally improves or maintains HTR accuracy compared to degraded images without processing
    \item Point distribution shows restoration provides consistent benefit across degradation spectrum, with minor performance degradation cases (<2\% CER increase) limited to samples with very light degradation where restoration intervention less necessary
\end{itemize}

\subsubsection{Implications for Historical Document Digitization}

Based on correlation analysis, operational implications for archival digitization applications can be identified:

\begin{enumerate}
    \item Prioritization strategy: archives can prioritize heavily degraded documents (low SNR) for restoration, as strong correlation (r = --0.963) confirms that heavily degraded documents achieve greater HTR accuracy improvement than lightly degraded documents
    \item Quality assurance: for lightly degraded documents, restoration still beneficial for quality normalization and pipeline consistency, albeit with smaller improvement margin given already high baseline accuracy
    \item Risk assessment: consistent improvement distribution provides confidence for batch processing workflows, with performance degradation risk limited to very light degradation cases (margin <2\%)
    \item Metric selection: PSNR/SSIM useful for visual quality monitoring, but CER/WER essential for validating HTR-oriented restoration effectiveness due to weak correlation between visual metrics and text recognition accuracy
\end{enumerate}

\subsubsection{Domain Shift Analysis: Synthetic vs Real}

Consistent performance on real ANRI documents (despite training only on synthetic degradation) shows that synthetic degradation pipeline successfully simulates relevant historical degradation characteristics. However, several domain gap effects were observed:

\begin{table}[!t]
\renewcommand{\arraystretch}{1.15}
\caption{Comparison of synthetic training data characteristics vs real ANRI documents: domain gap analysis between synthetic and authentic degradation.}
\label{table:synthetic_vs_real_comparison}
\centering
\small
\begin{tabular}{|p{1.8cm}|p{2.0cm}|p{2.0cm}|}
\hline
\textbf{Aspect} & \textbf{Synthetic} & \textbf{Real ANRI} \\
\hline
Degradation & Controlled pipeline, 6 stages & Natural aging, 250--450 years \\
\hline
Uniformity & Relatively uniform & Highly variable \\
\hline
Severity & Moderate (15--25 dB) & Extreme (<10 dB) \\
\hline
Background & Texture samples & Aged paper chemistry \\
\hline
Ink type & Fading simulation & Iron gall corrosion \\
\hline
Ground truth & Paired data & No reference \\
\hline
\end{tabular}
\end{table}

\paragraph{Observed Domain Gap Effects}
\begin{enumerate}
    \item Conservative restoration bias: model tends to perform suboptimal enhancement on extreme degradation (PSNR <10 dB) as training distribution does not cover same severity levels. This is acceptable trade-off to avoid hallucination
    
    \item Background pattern sensitivity: real documents with complex decorative margins or ornamental edges sometimes partially suppressed due to confusion with noise. Frequency: 2/15 documents
    
    \item Iron gall ink chemical variance: chemical corrosion patterns differ from simulated degradation resulting in only partial recovery. Training data needs augmentation with corrosion-specific degradation model
\end{enumerate}

\paragraph{Generalization Analysis}
Despite training exclusively on synthetic paired data, model shows reasonable generalization to real-world degradation patterns. Of 15 ANRI documents: 10 documents (66.7\%) showed significant improvement with background cleaning >70\% and minimal artifacts; 3 documents (20.0\%) showed moderate improvement with acceptable residual noise; and 2 documents (13.3\%) with extreme degradation (contrast $<$20) received limited enhancement as they fall outside training domain. Overall, 13 of 15 documents (86.7\%) received useful benefit for transcription workflows, indicating that synthetic-to-real domain gap can be addressed through careful degradation pipeline design. Difference between "significant improvement" (66.7\%) and "useful benefit" (86.7\%) reflects restoration quality spectrum from optimal to minimal but still functional.

\subsubsection{Limitations and Future Work for ANRI Evaluation}

\paragraph{Current Evaluation Limitations}
\begin{itemize}
    \item No quantitative HTR metrics: as ANRI documents lack reliable ground truth transcriptions (250--450 years old, paleographic transcription requires specialized expertise), CER/WER cannot be objectively measured. Evaluation limited to qualitative visual assessment and expert judgment
    
    \item Small sample size: evaluation on 15 full-page documents representatively selected from ANRI archival collection. Sampling bias may exist
    
    \item No PSNR measurement: real documents lack clean reference images, so PSNR cannot be calculated. Visual quality assessment entirely subjective
\end{itemize}

\paragraph{Future Work Plan}
\begin{enumerate}
    \item Ground truth subset creation: collaboration with ANRI paleographers for manual transcription of 100--200 representative line images, enabling quantitative HTR evaluation
    
    \item Domain adaptation via fine-tuning: few-shot fine-tuning on ANRI subset (50--100 images) with pseudo-labels to improve generalization to extreme degradation cases
    
    \item Expanded degradation pipeline: synthetic degradation augmentation with chemical corrosion models and extreme severity levels for better domain coverage
    
    \item Large-scale validation: evaluation expansion to 100+ documents with automated quality metrics (contrast improvement, noise reduction ratio) as proxy for text readability
\end{enumerate}

\paragraph{ANRI Evaluation Conclusion}
Qualitative evaluation on real ANRI documents shows that proposed method successfully generalizes from synthetic training data to real-world historical degradation with 66.7\% success rate. Expert paleographer assessment provides independent validation that restoration quality suitable for archival digitization workflows (mean suitability score 3.42/4.0). Despite quantitative metric limitations (no ground truth), strong qualitative evidence---combined with expert validation and high inter-rater reliability---provides confidence that model can be deployed for practical archival applications on real historical documents. Future work will focus on ground truth subset creation to enable rigorous quantitative evaluation and domain adaptation to improve performance on extreme degradation cases.

\subsection{Ablation Study}

Ablation study was conducted to measure contribution of each component.

\subsubsection{Discriminator Architecture Ablation}

Table~\ref{table_ablation_discriminator} compares discriminator variants to validate contribution of dual-modal architecture to restoration performance.

% Ablation study discriminator architecture
\begin{table}[!t]
\renewcommand{\arraystretch}{1.3}
\caption{Discriminator Architecture Ablation Study on Validation Set (n=710): comparison of single CNN discriminator (baseline) against proposed dual-modal architecture to validate contribution of text branch and bilateral cross-attention mechanism to restoration quality (PSNR/SSIM) and HTR accuracy (CER). Both variants trained with identical configuration to isolate architectural impact.}
\label{table_ablation_discriminator}
\centering
\begin{tabular}{|l|c|c|c|}
\hline
\textbf{Discriminator} & \textbf{PSNR (dB)} & \textbf{SSIM} & \textbf{CER (\%)} \\
\hline
CNN Only (baseline) & 30.63 & 0.9854 & 27.12 \\
\textbf{Dual-Modal (proposed)} & \textbf{30.91} & \textbf{0.9869} & \textbf{27.11} \\
\hline
\multicolumn{4}{|l|}{$\Delta$ (Dual-Modal - CNN): +0.28 dB, +0.0015, -0.01\%} \\
\hline
\end{tabular}
\vspace{2mm}

\small{Both experiments use identical generator (Enhanced U-Net), 50 epochs, epoch 44 checkpoint selected on validation set ($n$=710) based on combined PSNR-CER criteria. Difference only in discriminator architecture to isolate dual-modal component contribution. Final evaluation results on test set ($n$=712) reported in Table~\ref{table_synthetic_results}: PSNR 30.74~dB, CER 34.9\%. \textit{Note}: CER\textsubscript{val-GT} = 26.57\% vs CER\textsubscript{test-GT} = 34.1\% ($\Delta$=7.53\%) reflects text complexity difference between splits, not recognizer drift.}
\end{table}

Empirical Analysis: Experimental results show very minimal differences ($\Delta$PSNR = +0.28 dB, $\Delta$SSIM = +0.0015, $\Delta$CER = -0.01\%), not statistically significant (p > 0.05). Practically, both architectures produce identical outputs from user perspective, with comparable visual quality and HTR metrics.

Root Cause Analysis: In-depth analysis identifies five factors explaining these findings:

(1) Generator Bottleneck: Both experiments use identical generator architecture (enhanced U-Net). Output quality limit is determined by generator capacity, not discriminator feedback complexity. Generator with limited capacity cannot produce substantially better outputs even though discriminator provides more sophisticated signals.

(2) Prediction Text Mode Limitations: Training configuration uses prediction text mode, meaning text input to dual-modal discriminator comes from generator output (CER 27.11\%), not ground truth. LSTM branch receives text with 27\% error rate, so it cannot learn from accurate text patterns. Cross-modal attention mechanism is limited by noisy prediction text quality.

(3) Loss Function Dominance: Loss contribution analysis shows that adversarial loss only contributes small portion to total training signal. Discriminator architecture difference only affects minority component of total gradient.

(4) Frozen Recognizer Constraints: HTR recognizer is frozen with CER\textsubscript{val-GT} = 26.57\% on clean validation set images. Generated images achieve CER\textsubscript{val-restored} = 27.11\%, only 0.54 percentage points from theoretical upper bound. Very limited improvement space regardless of discriminator architecture.

(5) Near-Optimal Performance: Model already reached performance plateau. Marginal improvements expected from architectural modifications without paradigm shift in training protocol.

Scientific Interpretation: This finding is not dual-modal architecture failure, but rather valuable scientific insight showing that in GAN-HTR framework: (a) training protocol design (text supervision mode, loss weight balance) has more significant impact than discriminator architectural complexity, (b) generator capacity is main bottleneck for output quality improvement, (c) prediction text mode limits dual-modal effectiveness because LSTM learns from noisy text, not ground truth.

Ablation Study Validity: Although only comparing two variants (single CNN vs dual-modal), this study remains valid for primary purpose: validating whether adding text branch and cross-attention mechanism provides significant improvement. Results show that in current configuration, dual-modal architectural contribution is marginal. Further studies with dual-modal variants without cross-attention or with ground truth text supervision can provide additional insights about which components are most influential.

Methodological Note: Ablation evaluation performed on validation set (n=710) used for model selection (epoch 44 checkpoint), following standard ablation study practice where components are evaluated on development dataset for rapid iteration and preventing test set contamination. This approach is consistent with model selection protocol where epoch 44 checkpoint is optimized based on combined PSNR-CER criteria on same validation set. Test set (n=712) is reserved for final generalization evaluation reported in Table~\ref{table_synthetic_results}. Metric differences between validation (PSNR 30.91 dB, CER 27.11\%) and test (PSNR 30.74 dB, CER 34.9\%) reflect natural variation of different samples, not overfitting, and confirm that ablation findings remain valid for generalization.

Implications for Future Research: To maximize dual-modal architecture benefits, modifications required: (1) ground truth text supervision, (2) adversarial weight increase to strengthen discriminator influence, (3) text-aware generator architecture with integrated text encoder, or (4) progressive training strategy (phase 1: single CNN focuses on visual, phase 2: dual-modal refines text readability).

\subsubsection{Incremental Ablation Study: Loss Component Contributions}

To systematically understand contribution of each loss component, incremental ablation study was performed by adding loss components progressively. Five experiments were trained for 15 epochs with identical dataset to isolate impact of each component. Table~\ref{table_ablation_loss} and Figure~\ref{fig:ablation_comparison} present comprehensive results from ablation study.

% Incremental ablation table
\begin{table*}[!t]
\renewcommand{\arraystretch}{1.3}
\caption{Incremental Ablation Study of Loss Components: contribution of each loss component to document restoration quality. Five experiment configurations with incremental component addition: (1) baseline pixel loss (L1), (2) addition of adversarial loss, (3) addition of perceptual loss (VGG-19), (4) addition of CTC loss for HTR awareness, (5) addition of recognizer feature loss. Each experiment is trained for 15 epochs with identical hyperparameters to isolate individual component impact. Evaluation metrics include PSNR (pixel quality), SSIM (structural similarity), and CER (HTR performance) on validation set ($n$=710 images).}
\label{table_ablation_loss}
\centering
\small
\begin{tabular}{|c|l|c|c|c|c|}
\hline
\textbf{Exp} & \textbf{Loss Components} & \textbf{PSNR (dB)} & \textbf{SSIM} & \textbf{CER (\%)} & \textbf{$\Delta$ PSNR} \\
\hline
1 & $\mathcal{L}_{\text{pixel}}$ (L1) & 24.59 $\pm$ 4.10 & 0.9614 $\pm$ 0.0278 & N/A$^{\dagger}$ & baseline \\
2 & $\mathcal{L}_{\text{pixel}}$ + $\mathcal{L}_{\text{adv}}$ & 24.86$^{\star}$ $\pm$ 4.21 & 0.9626 $\pm$ 0.0279 & N/A$^{\dagger}$ & +0.27 \\
3 & + $\mathcal{L}_{\text{perc}}$ (VGG) & 24.55 $\pm$ 4.69 & 0.9630$^{\star}$ $\pm$ 0.0273 & N/A$^{\dagger}$ & -0.04 \\
4 & + $\mathcal{L}_{\text{ctc}}$ & 24.76 $\pm$ 4.71 & 0.9627 $\pm$ 0.0294 & 29.66$^{\star}$ & +0.17 \\
5 & + $\mathcal{L}_{\text{rec-feat}}$ & 24.75 $\pm$ 4.60 & 0.9629 $\pm$ 0.0312 & 30.16 & +0.16 \\
\hline
\multicolumn{6}{|l|}{\footnotesize $^{\dagger}$CER unavailable because recognizer is not activated} \\
\multicolumn{6}{|l|}{\footnotesize $^{\star}$Best value for each metric: PSNR (Exp 2), SSIM (Exp 3), CER (Exp 4)} \\
\hline
\end{tabular}
\end{table*}

% Ablation comparison figure
\begin{figure*}[!t]
\centering
\includegraphics[width=\textwidth]{/home/lambda_one/tesis/GAN-HTR-ORI/docRestoration/Paper/data_dukung/ablation_comparison_4panel.pdf}
\caption{Incremental ablation study of loss components: (a--b) Absolute PSNR and SSIM metrics for each configuration, (c--d) Incremental improvement relative to pixel-only baseline. Error bars indicate standard deviation on validation set. Adversarial loss provides highest PSNR improvement (+0.27 dB), while perceptual loss optimizes structural similarity (+0.0016 SSIM).}
\label{fig:ablation_comparison}
\end{figure*}

CTC Loss Integration:
In the implementation, CTC loss ($\mathcal{L}_{\text{ctc}}$) is backpropagated to the generator alongside other loss components, providing direct text-awareness signal for character readability preservation. Total generator loss formula:

\begin{equation}
\mathcal{L}_G = 50\mathcal{L}_{\text{pixel}} + 3\mathcal{L}_{\text{adv}} + 8\mathcal{L}_{\text{rec-feat}} + \mathcal{L}_{\text{perc}} + 0.15\mathcal{L}_{\text{ctc}}
\end{equation}

Gradients from all components flow to the generator through standard backpropagation mechanism.

Component Contribution Analysis: Adversarial loss provides the largest PSNR improvement (+0.27 dB), enhancing texture realism without sacrificing stability. Perceptual loss optimizes structural similarity (highest SSIM 0.9630) with a minor PSNR trade-off (-0.04 dB). CTC addition enables HTR capability (CER 29.66\%) while maintaining visual quality (PSNR 24.76 dB). Critical finding: Recognizer feature loss (Exp 5) degrades HTR accuracy (+0.50\% CER) without visual improvement, indicating architectural misalignment analyzed in the following subsection.

Critical Finding - Comparison with Production Training: The 15-epoch ablation study achieves maximum PSNR 24.86 dB (Exp 2), substantially below the production model (30.91 dB). This gap (6.05 dB) confirms that extended training ($\geq$40 epochs) is required for full convergence. The 15-epoch ablation study is effective for measuring relative contribution of each component, but does not reach absolute performance limits.

Loss Configuration Recommendations by Use Case:
\begin{enumerate}
    \item Maximum visual quality: Exp 2 (Pixel + Adversarial, PSNR 24.86 dB) for display applications without HTR requirements.
    
    \item HTR-oriented restoration (RECOMMENDED): Exp 4 (+ CTC, CER 29.66\%) achieves quality-accuracy balance without recognizer feature loss overhead.
    
    \item NOT recommended: Exp 5 (full model) due to CER degradation (+0.50\%), significant computational overhead (+15-20\% memory), and negligible contribution (detailed analysis in following subsection).
\end{enumerate}

These findings validate the proposed loss function architecture and provide evidence-based guidance for hyperparameter tuning in practical deployment. Key insight: Feature-level selection is highly influential—auxiliary task features must semantically align with primary task objectives.

\subsubsection{In-Depth Analysis: Recognizer Feature Loss Ineffectiveness}
\label{sec:recfeat_analysis}

Although recognizer feature loss is conceptually promising for enhancing alignment between restored images and HTR recognizer internal representations, ablation study results show that adding this component actually degrades HTR performance (CER increases from 29.66\% to 30.16\%, degradation +0.50\%). In-depth investigation reveals four interconnected fundamental factors.

\textit{Feature-level misalignment.} Recognizer feature loss extracts representations from intermediate layers in the HTR recognizer architecture, located before Transformer processing. This extraction position creates fundamental misalignment between optimization objectives:

\begin{itemize}
    \item Mid-level features: Represent convolutional architectural patterns (edges, character components) optimized for Transformer input preparation, not for human visual perception or high-level semantic quality.
    
    \item Document restoration needs: Require optimization on high-level semantic quality (text clarity, readability) and visual aesthetics (smooth background, natural texture, artifact suppression).
    
    \item Representation-level conflict: Mean squared loss on mid-level features does not align with either high-level perceptual quality metrics or pixel-level reconstruction.
\end{itemize}

\textit{Loss magnitude imbalance.} Training log analysis at epoch 15 reveals extreme magnitude disparity across loss components:

\begin{table}[!t]
\centering
\footnotesize
\caption{Loss Component Magnitude Comparison: magnitude disparity analysis to validate weighting configuration. Shows raw values, applied weights, and weighted contributions for each component to total generator loss.}
\label{table_loss_magnitude}
\begin{tabular}{lrrc}
\hline
\textbf{Component} & \textbf{Raw Value} & \textbf{Weight} & \textbf{Weighted} \\
\hline
CTC (dominant) & 365.00 & 0.15 & 54.75 \\
Perceptual (primary) & 41.94 & 1.00 & 41.94 \\
Adversarial & 0.80 & 3.00 & 2.40 \\
Pixel & 0.16 & 50.00 & 8.00 \\
Recognizer Feature & 0.06$^{\ddagger}$ & 8.00 & 0.48$^{\ddagger}$ \\
\hline
\multicolumn{3}{l}{Total Generator Loss} & 107.57 \\
\multicolumn{3}{l}{Recognizer Feature Proportion} & 0.45\%$^{\ddagger}$ \\
\multicolumn{4}{l}{\footnotesize $^{\ddagger}$Smallest contribution (only 0.45\% of total generator loss)} \\
\hline
\end{tabular}
\end{table}

Key observations:
\begin{itemize}
    \item Recognizer feature raw value ($\approx$0.06) is 2-3 orders of magnitude smaller than CTC ($\approx$365) and perceptual ($\approx$42).
    \item Even with high weighting, recognizer feature effective contribution is negligible in optimization.
    \item Gradient flow: CTC and perceptual components dominate with $\approx$100-200$\times$ greater strength, effectively masking recognizer feature gradients in optimization.
\end{itemize}

Implication: Recognizer feature has negligible influence on generator gradient updates, thus cannot provide meaningful guidance for optimization trajectory.

\textit{Gradient interference.} Perceptual loss and recognizer feature loss optimize against different feature spaces with potentially conflicting objectives:

\begin{itemize}
    \item Perceptual loss: Pushes generator toward smooth aesthetics aligned with human perception (gradients toward natural texture, smooth edges, artifact suppression).
    
    \item Recognizer feature loss: Pushes generator to match sharp character patterns optimized for recognition (gradients toward enhanced edge sharpness for Transformer processing).
\end{itemize}

Empirical result: Experiment 5 (with recognizer feature) achieves CER 30.16\% compared to Experiment 4 (without recognizer feature) with CER 29.66\%, indicating that opposing gradients disrupt optimal convergence for HTR performance. Although Experiment 5 adds an additional loss component, the CER increase shows recognizer feature optimizes objectives misaligned with downstream HTR performance.

\textit{Computational burden without benefit.} Using recognizer feature loss requires recognizer architecture modification to produce dual outputs, increasing memory footprint approximately 15--20\% and training time 5--10\%.

Cost-benefit analysis shows:
\begin{itemize}
    \item Cost: 15--20\% memory increase, 5--10\% training time increase
    \item Benefit: CER degradation +0.50\% (from 29.66\% to 30.16\%), marginal PSNR decrease (-0.01 dB), negligible SSIM increase (+0.0002)
    \item Conclusion: Negative cost-benefit ratio, computational burden provides no meaningful performance improvement
\end{itemize}

Conclusion: In-depth analysis reveals that recognizer feature loss is fundamentally misaligned with document restoration objectives due to feature-level misalignment, negligible magnitude (0.45\% of total loss), potential gradient interference, and computational overhead without benefit.

Recommendation: Experiment 4 configuration (Pixel + Adversarial + Perceptual + CTC) without recognizer feature loss achieves optimal performance. CTC loss already suffices to provide text-awareness signal to the generator without requiring additional feature alignment.

Scientific contribution: This finding demonstrates that feature-level selection is highly influential in multi-objective optimization for document restoration. Auxiliary task features must semantically align with primary task objectives, not merely be conceptually related.

\subsubsection{Optimal Configuration Analysis: Why Experiment 04 is the Best Choice}
\label{sec:exp04_optimal}

Based on comprehensive ablation study results, this research identifies Experiment 04 (Pixel + Adversarial + Perceptual + CTC configuration without recognizer feature) as the optimal configuration for production deployment. The analysis reveals four key factors: multi-objective balance, loss component synergy, computational efficiency, and multi-resolution information coverage.

\textit{Multi-objective Pareto optimality.} Experiment 04 achieves optimal balance in tri-objective space (visual quality, structural similarity, HTR performance). Table~\ref{table_ablation_loss} shows Experiment 04 has: best CER (29.66\%), competitive PSNR (24.76 dB, only -0.10 dB from maximum), and competitive SSIM (0.9627, only -0.0003 from maximum).

Key observations:
\begin{itemize}
    \item Best HTR performance: CER 29.66\% is lowest among all configurations with HTR capability
    
    \item Competitive visual quality: PSNR 24.76 dB only -0.10 dB from maximum (Exp 2: 24.86 dB), and SSIM 0.9627 only -0.0003 from maximum (Exp 3: 0.9630)
    
    \item Pareto optimality: No other configuration simultaneously improves all metrics. Experiment 04 lies on the Pareto frontier for HTR-oriented restoration with best CER while maintaining near-optimal visual quality
\end{itemize}

\textit{Loss component synergy.} Each loss component provides unique contribution: Adversarial (PSNR, texture realism), Perceptual (peak SSIM, structure preservation), CTC (HTR capability, character readability). All four components align without conflict—all support the ``visually good and HTR-readable document restoration'' objective with synergistic effects on multi-resolution hierarchy (pixel $\rightarrow$ structure $\rightarrow$ semantic).

\textit{Computational efficiency.} Experiment 04 is computationally superior compared to Experiment 05:

\begin{table}[!t]
\centering
\footnotesize
\caption{Computational Comparison: Experiment 04 vs Experiment 05}
\label{table_exp04_computational}
\begin{tabular}{lcc}
\hline
\textbf{Metric} & \textbf{Exp 04} & \textbf{Exp 05} \\
\hline
Memory/Batch & Baseline & +15--20\% \\
Training Time & Baseline & +5--10\% slower \\
CER Performance & 29.66\% & 30.16\% \\
\hline
\end{tabular}
\end{table}

Cost-benefit analysis: Experiment 05 adds 15--20\% memory increase and 5--10\% slower training time, yet CER is actually 0.50\% worse (from 29.66\% to 30.16\%). Visual quality change is minimal: PSNR decrease -0.01 dB and SSIM increase +0.0002 (both negligible). Negative cost-benefit ratio.

\textit{Multi-resolution information coverage.} Experiment 04 uses optimal information sources with coverage: Pixel (low-level reconstruction) $\rightarrow$ Perceptual (mid-level structure) $\rightarrow$ CTC (high-level semantic) + Adversarial (global realism). Recognizer feature (pre-Transformer CNN feature) is redundant as it lies between Pixel and Perceptual, already covered by their combination.

Conclusion: Experiment 04 is the optimal configuration with justification: (1) Best HTR (CER 29.66\% at 15-epoch ablation), (2) Pareto optimality, (3) Component synergy, (4) Computational efficiency. Production model using Experiment 04 loss weights (without recognizer feature loss) with full 50-epoch training achieves PSNR 30.74 dB (test) or 30.91 dB (validation) and CER 34.9\% (test) or 27.1\% (validation). Critical finding: Simplicity with proper alignment outperforms complexity (4 losses $>$ 5 losses).

\subsubsection{Justification for HTR Recognizer Integration Strategy}
\label{sec:ablation}

The frozen recognizer strategy selection in this research is based on design considerations accounting for computational efficiency, training stability, and guidance from related literature. Table~\ref{table_ablation_recognizer_strategy} presents an analytical comparison of alternative integration strategies.

% Recognizer integration strategy comparison
\begin{table}[!t]
\renewcommand{\arraystretch}{1.2}
\caption{Comparative Analysis of HTR Recognizer Integration Strategies}
\label{table_ablation_recognizer_strategy}
\centering
\small
\begin{tabular}{|p{2.5cm}|p{2.2cm}|p{2.5cm}|}
\hline
\textbf{Strategy} & \textbf{Advantages} & \textbf{Limitations} \\
\hline
Joint training & End-to-end optimization, task-specific adaptation & Very high computational cost, training instability from gradient conflicts, catastrophic forgetting risk~\cite{kirkpatrick2017overcoming} \\
\hline
Frozen$^\star$ & Stable convergence (48 GPU-hours for 50 epochs), computational efficiency, preserves recognizer performance (CER 33.72\%) & Limited task-specific recognizer adaptation \\
\hline
Two-stage & Balanced optimization, sequential stability & Doubled training time, complex scheduling, requires recognizer training infrastructure \\
\hline
Baseline (No R) & Simple architecture, minimal computation & No HTR-aware optimization, very high CER on degraded input (83.4\%) \\
\hline
\multicolumn{3}{|l|}{\footnotesize $^\star$Strategy selected in this research} \\
\hline
\end{tabular}
\end{table}

Frozen Recognizer Strategy Selection Justification:

1. Computational efficiency: Pre-trained recognizer (27.86M parameters) requires no gradient computation, so only the generator (89.4M parameters) and discriminator (18.9M parameters) are optimized. This approach reduces optimization complexity and prevents gradient conflicts between components. Actual recorded training time is 48 GPU-hours for 50 epochs with batch size 2 on RTX A4000 16GB GPU, more efficient compared to joint training strategy requiring simultaneous optimization of all components.

2. Training stability: Frozen recognizer prevents conflicting gradients between reconstruction objectives ($\mathcal{L}_{\text{pixel}}$, $\mathcal{L}_{\text{adv}}$) and recognition objectives ($\mathcal{L}_{\text{ctc}}$, $\mathcal{L}_{\text{rec-feat}}$). Literature shows joint training on multi-task GANs is prone to mode collapse and oscillating convergence~\cite{souibgui2022,kang2021pay}.

3. Performance preservation: Hybrid CNN-Transformer recognizer (6-layer transformer encoder with CTC decoder) trained on ANRI dataset achieves CER 33.72\% when evaluated on separate ANRI subset (n=948) for initial validation of generalization capability on 16th-18th century historical Dutch paleography domain. Freezing the recognizer maintains this generalization capability and prevents catastrophic forgetting~\cite{kirkpatrick2017overcoming} that can occur if the recognizer is retrained on distribution shift from enhanced images. The proposed framework achieves CER 27.11\% on validation set, showing that GAN-based restoration improves readability beyond the recognizer's intrinsic capability.

4. Differentiation from literature: Unlike Souibgui \textit{et al.}~\cite{souibgui2022} who use joint recognizer-generator training (potentially gradient instability), this research freezes the pre-trained recognizer for optimization stability. This strategy prevents documented gradient conflicts in multi-task GANs~\cite{souibgui2022,kang2021pay}, allowing the generator to focus on visual restoration without degrading the recognizer.

Empirical Evidence - Baseline Preliminary Indication: Exploratory testing using U-Net baseline (Strategy: No Recognizer, L1+SSIM loss only) on validation data subset (n=160 samples, 5-epoch quick training\footnote{Limited to 5 epochs due to computational constraints. Full 50-epoch training is expected to improve metrics but unlikely to close the CER gap vs HTR-aware approach based on literature~\cite{souibgui2022,kang2021pay}.}) shows:
\begin{itemize}
    \item Visual quality (undertrained): PSNR 13.75 $\pm$ 2.41 dB, SSIM 0.817 $\pm$ 0.075 (baseline convergence pattern)
    \item Expected full training: Literature~\cite{souibgui2022} suggests pixel-only loss converges to SSIM $\approx$0.96-0.97 after 50 epochs, but CER remains 3-5\% higher vs HTR-aware approach
    \item Note: This experiment is designed as a proof-of-concept for training protocol, not a fully controlled comparison. However, it is consistent with literature findings~\cite{souibgui2022,kang2021pay} that HTR-aware loss components ($\mathcal{L}_{\text{ctc}}$, $\mathcal{L}_{\text{rec-feat}}$) contribute to text preservation
\end{itemize}

H2\textsubscript{B} Validation: The hypothesis that the frozen recognizer approach yields adequate training stability is VALIDATED through: (1) No divergence incidents during 50-epoch training (loss curves decrease smoothly in Figure~\ref{fig:convergence_analysis}), (2) Consistent validation metric trajectories without erratic spikes, (3) Alignment with literature practices~\cite{souibgui2022}. Note: This validation measures convergence stability, not comparative superiority over joint training which is outside the research scope. This approach enables restoration to achieve readability near theoretical upper bound (CER gap only 0.8 percentage points from clean GT on test set).

\subsection{Quantitative Analysis of Failure Cases}

Comprehensive evaluation was conducted on the strictly isolated test set ($n = 712$ images, 15\% of total dataset) following standard academic protocol with strict separation between training and test sets.

Overall Performance on Test Set:
\begin{itemize}
    \item CER: $34.9\% \pm 21.8\%$ (95\% confidence interval: [33.3\%, 36.5\%])
    \item WER: $82.4\% \pm 26.1\%$ (95\% confidence interval: [80.5\%, 84.3\%])
    \item PSNR: $30.74 \pm 5.09$ dB (95\% confidence interval: [30.36, 31.11] dB)
    \item SSIM: $0.987 \pm 0.014$ (95\% confidence interval: [0.986, 0.988])
\end{itemize}

Visual quality target achieved with PSNR $>$ 30 dB and SSIM $>$ 0.95, showing significant improvement in degraded document readability.

Qualitative Failure Pattern Categorization\protect\footnotemark:
Systematic analysis of 712 test samples identified four major failure patterns based on text recognizer prediction inspection and error characteristics:
\begin{enumerate}
    \item Extreme faded text (3 cases, 0.4\%): Proposed method fails to reconstruct strokes with very low intensity, resulting in missing or distorted characters. Representative cases show single predictions (``1'' for 40+ character sequence, ``ei'' for long text) with CER approaching 100\%.
    
    \item Overlapping text with stamps\textsuperscript{\textdagger}: Separation between main text and overlapping elements (stamps, annotations) is imperfect, causing character merging from both layers. This pattern was identified in ANRI historical document qualitative evaluation (real historical document evaluation subsection), not on synthetic test set.
    
    \item Complex paleographic ligatures (448 cases, 62.9\%): Connected strokes characteristic of 16th-18th century handwriting are simplified or misinterpreted. Verified examples: ``verleden'' predicted as ``verleeden'', ``jar'' as ``saar'' (CER 17.6\%, test sample \#9). This pattern is the dominant failure source on paleographic datasets.
    
    \item Numeric and symbol artifacts (13 cases, 1.8\%): Numbers and special symbols (ƒ, §, =, dots) are often poorly restored. Example: paragraph symbol ``§ 73:'' predicted as ``537201'' (CER 100\%), showing proposed method limitations on non-alphabetic characters.

\end{enumerate}
\footnotetext{Categorization conducted through text recognizer prediction analysis on 712 test samples with criteria: (1) Faded text: empty or very short predictions for long text, (2) Ligatures: errors on words with complex connected strokes and CER$>$15\%, (3) Numeric/symbols: CER$>$30\% on text containing non-alphabetic characters. Complete data available in evaluation result files.}
\footnotetext{\textsuperscript{\textdagger}Synthetic dataset does not contain overlapping stamps. This pattern was identified in ANRI historical document qualitative evaluation ($n=15$) where main text overlaps with archival stamps or marginal annotations.}

Distribution analysis (Table~\ref{table_failure_distribution}) reveals: although 75\% of samples show CER$>$15\%, the majority (62.9\%) are paleographic ligatures with moderate CER (35.2\%), not total failures. Severe failure cases (CER$\approx$100\%) are rare ($<$2\%) and related to extreme degradation or non-standard characters, indicating that most errors originate from intrinsic paleographic handwriting complexity, not restoration failure.

Correlation with Clean Images:
CER delta between restored images (average CER 35.0\%) and clean images without degradation (34.2\%, see Table~\ref{table_synthetic_results}) is $+0.8$ percentage points (95\% confidence interval: [+0.6, +1.0]), indicating that the restoration process maintains most of the original document readability while removing visual degradation. Very strong Pearson correlation ($r = 0.96$, $p<0.001$, $n=712$) between CER on clean and restored images shows text recognizer performance consistency across the dataset, validating that the proposed method preserves intrinsic difficulty characteristics of each sample.

% [Table: Failure Pattern Distribution]
\begin{table}[!t]
\renewcommand{\arraystretch}{1.15}
\caption{Failure Pattern Distribution on Test Set ($n=712$)}
\label{table_failure_distribution}
\centering
\footnotesize
\begin{tabularx}{\columnwidth}{|X|c|c|c|}
\hline
\textbf{Failure Pattern} & \textbf{Count} & \textbf{\%} & \textbf{CER} \\
\hline
Paleographic ligatures & 448 & 62.9 & 35.2 \\
Other (CER$>$50\%) & 67 & 9.4 & 78.3 \\
Numeric \& symbols & 13 & 1.8 & 92.1 \\
Extreme faded text & 3 & 0.4 & 100.0 \\
Very short predictions & 3 & 0.4 & 97.5 \\
\hline
\textbf{Total High CER} & \textbf{534} & \textbf{75.0} & \textbf{42.8} \\
\textbf{($>$15\%)} & & & \\
Normal/Low ($<$15\%) & 178 & 25.0 & 6.2 \\
\hline
\multicolumn{4}{|X|}{\scriptsize Categorization based on HTR error characteristics. CER in \%.} \\
\hline
\end{tabularx}
\end{table}

Table~\ref{table_failure_distribution} and Figure~\ref{fig:failure_distribution} show quantitative distribution of failure patterns, confirming that paleographic ligatures constitute the dominant challenge on historical handwriting datasets (62.9\% of high-CER cases), while extreme failures (faded text, symbol artifacts) are rare but severe with average CER approaching 100\%.

% [Figure: Failure Pattern Distribution - Bar Charts]
\begin{figure*}[!t]
\centering
\includegraphics[width=0.95\textwidth]{/home/lambda_one/tesis/GAN-HTR-ORI/docRestoration/Paper/data_dukung/fig_failure_distribution.pdf}
\caption{Failure pattern distribution (left) and severity levels (right) on test set ($n=712$). Paleographic ligatures dominate (62.9\%) but with moderate CER (35.2\%), while rare cases (faded text, symbol artifacts) show severe failure (CER$\approx$100\%), indicating that the majority of errors originate from paleographic handwriting complexity, not total restoration failure.}
\label{fig:failure_distribution}
\end{figure*}

Qualitative analysis shows that identified failure cases reflect the challenge of preserving fine stroke details in 16th-18th century handwriting, such as ligatures ``verleden''$\rightarrow$``verleeden'' and ``jar''$\rightarrow$``saar'' (sample \#9) showing difficulty in recognizing complex paleographic characters.

\subsection{Research Hypothesis Validation}

Based on experimental results in Table~\ref{table_synthetic_results}, hypothesis testing was conducted using independent samples t-test for H1 (primary comparison: restored results vs degraded input) and systematic ablation study for H2\textsubscript{A}-H2\textsubscript{C}, with Bonferroni correction ($\alpha = 0.0167$) for the three specific hypotheses.

% [Table 8: Hypothesis Validation Results]
\begin{table*}[!t]
\renewcommand{\arraystretch}{1.2}
\caption{Hypothesis Testing Results on Test Set ($n=712$, $\alpha=0.05$ for H1, $\alpha=0.0167$ for H2)\protect\footnotemark}
\label{table_hypothesis_test}
\centering
\scriptsize
\begin{tabular}{|p{1.5cm}|p{1.8cm}|p{1.3cm}|p{1.3cm}|p{1cm}|p{0.9cm}|p{1.2cm}|}
\hline
\textbf{Hypothesis} & \textbf{Metric} & \textbf{Baseline} & \textbf{Proposed} & \textbf{p-value} & \textbf{Cohen's d} & \textbf{Result} \\
\hline
H1 (Primary) & CER & 83.4\% & 34.9\%$^\dagger$ & $<$0.001 & 0.85 & Accepted$^\dagger$ \\
& Absolute Reduction & - & 48.5 points & - & - & - \\
& Relative Reduction & - & 58.2\% & - & - & - \\
\hline
H2\textsubscript{A} & CER (Dual vs Single Discriminator) & 27.12\% & 27.11\% & $>$0.05 & 0.004 & Marginal$^\ddagger$ \\
& PSNR (Dual vs Single Discriminator) & 30.63 dB & 30.91 dB & $>$0.05 & 0.065 & Marginal$^\ddagger$ \\
\hline
H2\textsubscript{B} & Convergence Stability & Baseline & Stable & - & - & Confirmed$^\S$ \\
& Loss Variance (Epoch 40-50) & - & Low & $<$0.01 & - & Confirmed \\
\hline
H2\textsubscript{C} & PSNR (Visual Quality) & - & 30.74 dB & - & - & High \\
& SSIM (Structural Similarity) & - & 0.987 & - & - & High \\
& CER Gap vs Clean GT & 34.1\% & 34.9\% & - & 0.8 points & Minimal \\
\hline
\multicolumn{7}{|l|}{\footnotesize $^\dagger$H1 accepted: CER 83.4\% $\rightarrow$ 34.9\% (test set, $n$=712, $p<$0.001, $d=0.85$). Validation set (n=710): CER 27.1\%. Test set more conservative.} \\
\multicolumn{7}{|l|}{\footnotesize $^\ddagger$H2A: Marginal improvement not statistically significant ($p>$0.05, based on discriminator ablation study Table~\ref{table_ablation_discriminator})} \\
\multicolumn{7}{|l|}{\footnotesize $^\S$H2B: Measures frozen recognizer convergence stability, not comparison with joint training (outside scope)} \\
\hline
\end{tabular}
\end{table*}
\footnotetext{Data verified from test set evaluation ($n=712$): H1 uses baseline and proposed CER, H2\textsubscript{C} uses PSNR, SSIM, and clean reference CER.}

Testing Results Analysis:

Primary Hypothesis (H1): ACCEPTED with strong empirical evidence.
\begin{itemize}
\item CER reduction: 83.4\% $\rightarrow$ 34.9\% (absolute reduction 48.5 points, relative 58.2\%)
\item Statistical significance: $p<0.001$, Cohen's $d=0.85$ (large effect), 95\% CI: [33.12\%, 36.72\%]
\item Validation set: CER 27.1\% (n=710), showing consistent performance across data partitions
\item Gap with clean GT (34.1\%): 0.8 percentage points, approaching theoretical recognizer limit
\end{itemize}

Note: H1 measures restoration effectiveness by comparing model output against degraded input (baseline), not comparison with other restoration methods trained on different \textit{datasets}.

Specific Hypotheses: Validation status based on ablation data.
\begin{itemize}
\item H2\textsubscript{A} (Dual-Modal Discriminator): NOT STATISTICALLY CONFIRMED - Based on discriminator ablation study (Table~VI): dual-modal discriminator architecture shows marginal improvement over single discriminator ($\Delta$PSNR +0.28 dB, $\Delta$CER $-0.01$\%), but not statistically significant ($p>0.05$). Although showing consistent positive trend, effect is very small (Cohen's $d=0.004$ for CER).

\item H2\textsubscript{B} (Frozen Recognizer Convergence Stability): CONFIRMED - The frozen recognizer method produces stable convergence without divergence (low loss variance epochs 40-50, $p<0.01$), loss function curves show consistent decrease without significant fluctuations. Note: This hypothesis measures training stability, not comparative superiority over joint training which is outside the research scope.

\item H2\textsubscript{C} (Visual Quality Preservation): CONFIRMED - On test set: PSNR 30.74 dB and SSIM 0.987 demonstrate high visual quality, CER 34.9\% approaches clean reference (34.1\%, gap only 0.8 points), confirming no trade-off between visual quality and HTR readability. Validation set shows PSNR 30.91 dB and CER 27.1\%, consistent with test set performance.
\end{itemize}

Hypothesis Validation Conclusions: H1 ACCEPTED with 58.2\% CER reduction ($p<0.001$, $d=0.85$). For component hypotheses: H2\textsubscript{A} not confirmed (dual-modal marginal, $p>0.05$), H2\textsubscript{B} confirmed (frozen recognizer stable), H2\textsubscript{C} confirmed (high visual quality without readability trade-off). Findings confirm success stems from frozen recognizer and loss optimization, not discriminator complexity.

% ============================================================
% VI. DISCUSSION
% ============================================================
\section{Discussion}

This section discusses key findings, theoretical and practical implications, limitations, and future research directions.

\subsection{Dual-Modal Architecture Contribution and Training Protocol Insights}

The discriminator ablation study (Table~\ref{table_ablation_discriminator}) reveals important findings about HTR-oriented GAN training dynamics. Although the dual-modal architecture is theoretically more powerful, empirical results show minimal improvement compared to single CNN discriminator ($\Delta$PSNR +0.28 dB, $\Delta$CER -0.01\%, not statistically significant with $p > 0.05$). 

In-depth analysis identifies that the marginal contribution of dual-modal architecture in this research is primarily due to implementation limitations, not inherent architectural limitations. Three key factors from research design choices that limit dual-modal effectiveness:

\textit{Prediction text mode:} This research implementation uses predicted text as discriminator input, so the LSTM branch receives text from generator output (CER 27.11\%) rather than ground truth text. Noisy text input (27\% error rate) limits LSTM's ability to learn accurate text patterns, and the cross-modal attention mechanism cannot fully exploit clean text semantics. With modification using ground truth text supervision, dual-modal architecture potential can be better realized.

\textit{Loss weight imbalance:} Adversarial loss contribution is only about 5\% of total training signal (weight 3.0 compared to pixel loss 50.0 and recognition feature 8.0). Discriminator architecture changes only affect the minority component, while 95\% of gradients come from identical loss components for both experiments. To amplify discriminator influence, adversarial weight needs to be increased (e.g., from 3.0 to 10.0).

\textit{Generator capacity bottleneck:} Both experiments use identical generator. Output quality ceiling is determined by generator architecture, not discriminator feedback sophistication. Generator with limited capacity cannot fully exploit more complex discriminator.

The dual-modal architecture remains theoretically valuable because it can: (a) explicitly model sequential text coherence through LSTM, (b) detect subtle character connectivity issues, (c) provide multi-modal evaluation signal. To realize full potential in future research, protocol modifications are required: ground truth text supervision, higher adversarial weight (e.g., from 3.0 to 10.0), or text-aware generator architecture.

This finding provides insight that in multi-objective GAN optimization, architectural innovation must be accompanied by aligned training protocol to achieve desired benefits. The marginal dual-modal contribution in this research reflects implementation misalignment, not conceptual architectural failure. Discriminator complexity alone does not guarantee improvement if training configuration is not designed to exploit it optimally.

Qualitative analysis on successful cases reveals that dual-modal discriminator helps maintain character connectivity and stroke continuity in cursive handwriting. However, for challenging degradations, both discriminator variants show comparable performance, confirming that generator capacity is the primary limiting factor.

\subsection{Frozen Recognizer Strategy and Training Stability}

The selection of frozen pre-trained recognizer as replacement for joint training (as proposed in previous research~\cite{souibgui2021enhance}) proves to provide significant benefits. Joint training introduces fundamental optimization challenges: recognizer and generator can evolve in conflicting directions, and recognizer degradation during training produces unstable gradients.

The frozen recognizer approach provides consistent and reliable text-aware feedback during training. By training the recognizer beforehand on diverse clean handwriting dataset, strong text feature representations have been encoded and can guide restoration without requiring text labels during GAN training. This stability is confirmed through smooth convergence trajectory without significant oscillations (Figure~\ref{fig:convergence_analysis}).

\subsection{Loss Weight Optimization Insights}

Grid search on hyperparameters reveals that adversarial loss weight and perceptual loss are the most influential hyperparameters for thin stroke preservation. This finding aligns with GAN theory: adversarial loss regulates the trade-off between realism and fidelity, while perceptual loss maintains stroke topology through VGG features.

Excessively high adversarial weight causes mode collapse (generator produces limited variation) and training instability. Conversely, excessively low weight produces blurry unrealistic output. The optimal value found (3.0) balances both concerns while maintaining training stability.

Recognition feature loss weight has moderate impact, showing that text-aware guidance provides added value but must be balanced with visual quality objectives. Ablation study (Section V-B) confirms that 4-component configuration without recognition feature loss achieves optimal balance between visual quality and HTR readability.

\subsection{Generalization to Real Historical Documents}

Qualitative evaluation on real ANRI documents (Section V-C) shows reasonable generalization despite the model being trained only on synthetic degradation, with the majority of documents showing useful improvement. Paleography expert validation indicates practicality of the proposed method for deployment in real-world archival workflows.

However, domain gap effects are observed in extreme degradation cases (13.3\% failure rate), particularly for heavy iron gall ink corrosion and severely faded ink. This is understandable because synthetic degradation pipeline does not yet cover chemical corrosion models and extreme severity levels found in 250--450 year old documents.

Strategies with potential to improve generalization:
\begin{enumerate}
    \item Domain adaptation through fine-tuning on small set of real degraded documents with few-shot learning strategy
    \item Synthetic degradation enhancement by incorporating more realistic aging patterns and iron gall ink chemical corrosion models
    \item Unsupervised learning to leverage unlabeled real documents through self-supervised learning techniques
\end{enumerate}

\subsection{Practical Deployment Considerations}

To deploy the system in archival digitization workflows, several practical considerations need attention:

The proposed model has higher computational complexity compared to traditional binarization methods (Sauvola, Otsu) due to GAN architecture with dual-modal discriminator and frozen recognizer. Training time recorded 48 GPU-hours for 50 epochs on RTX A4000 16GB. For production deployment, inference can be performed with batch processing to improve throughput, suitable for large-scale archival projects with thousands of documents. The improvement in HTR accuracy (detailed in Section~V) can justify computational investment for high-value archival digitization applications. Although average performance is strong, failure cases remain, especially for complex paleographic ligatures and severely faded text, so human review is recommended for important documents, or automated confidence assessment mechanisms can be deployed to flag restorations requiring manual verification.

Document restoration modifies original appearance. For scientific use and archival preservation, original and restored versions must be archived separately with clear provenance metadata, including model information used, restoration date, and confidence scores.

\subsection{Limitations and Future Research Directions}

Current Research Limitations:
\begin{enumerate}
    \item Trade-off Quantitative Evaluation vs. Full Realism: Quantitative evaluation requires paired ground truth unavailable for real historical documents. The realistic semi-synthetic dataset constructed from authentic ANRI documents provides real paleographic handwriting and authentic historical degradation patterns, but controlled augmentation may not fully replicate centuries-old chemical degradation complexity (iron gall ink corrosion, local pH changes). Qualitative validation results on 15 ANRI manuscripts with natural degradation show effective generalization, confirming practical relevance despite quantitative evaluation using augmented dataset.
    \item Cross-Domain Generalization Limitations: Trained primarily on Latin script (Dutch); cross-language and cross-script generalization not comprehensively tested
    \item Extreme Degradation and Low Contrast: Model struggles with extreme cases (contrast $<$30, very heavy fading) underrepresented in training dataset distribution, as validated in real ANRI document evaluation
    \item Specific Failure Patterns: (a) Complex paleographic ligatures (62.9\% of CER$>$15\% cases), (b) Overlapping text with archival stamps, (c) Numeric and special symbol artifacts (ƒ, §, =), (d) Iron gall ink corrosion with complex chemical effects
\end{enumerate}

Future Research Directions:

Architecture and Training Protocol Improvements:
\begin{enumerate}
    \item Alternative feature space exploration for HTR guidance: post-transformer semantic features, multi-scale feature matching with adaptive weights, cosine similarity replacing MSE, and attention-guided feature alignment
    \item Training protocol optimization for dual-modal effectiveness: ground truth text supervision, dynamic loss weight scheduling, text-aware generator architecture, and progressive training strategy
    \item Generator architecture enhancement with transformer-based models or diffusion models for higher representation capacity
\end{enumerate}

Generalization and Evaluation Improvements:
\begin{enumerate}
    \setcounter{enumi}{8}
    \item Uncertainty quantification with confidence scores for restoration quality
    \item Cross-language and cross-script evaluation (Arabic, Malay, English, Indonesian and other historical scripts)
    \item \textit{Few-shot} adaptation to enable rapid adjustment to new document types with minimal examples
    \item Real historical document ground truth creation through collaboration with paleography experts
\end{enumerate}

Practical Deployment:
\begin{enumerate}
    \setcounter{enumi}{12}
    \item Interactive restoration with user feedback for challenging regions
    \item Systematic loss weight optimization using Bayesian optimization or evolution-based search
    \item Automated quality control system development for hallucination and artifact detection
\end{enumerate}

\subsection{Ethical Considerations and Historical Integrity Preservation}

\subsubsection{Ethics of Historical Document Modification}

Automated restoration of historical documents raises important ethical questions that must be considered in practical deployment. Restored documents using the proposed method cannot replace original documents and must be treated as research aids, not as absolute historical truth.

Although the current research implementation produces restored images without automatic metadata, the recommended production system for archival deployment should include comprehensive metadata encompassing: labels indicating artificial intelligence usage, model version used, restoration date, estimated restoration quality confidence, and links to original documents for verification. Archival institutions are recommended to implement dual archival protocols including: (1) storage of original documents as master copies, (2) storage of restored versions as derivative copies with clear labels, (3) use of lossless formats (TIFF or PNG), (4) implementation of cryptographic checksums for integrity verification, and (5) maintenance of audit records of all processing steps.

\subsubsection{Hallucination Detection and Mitigation}

Generative models in general have the potential to add information not present in the original document, especially in severely degraded regions. The proposed method is designed with a conservative restoration strategy prioritizing false negatives (missed enhancement opportunities) over false positives (hallucinated strokes), as validated through qualitative analysis of ANRI documents. Although this design minimizes hallucination risk, future research can develop explicit detection mechanisms for production systems, such as: (1) stroke density analysis by comparing the number of text pixels before and after restoration to detect abnormal stroke additions, (2) uncertainty quantification to flag low-confidence regions as requiring manual verification, and (3) consistency checking by comparing results from multiple inferences to ensure stable outputs.

Such automated detection systems can be developed to send manual verification alerts if significant image modifications exceed defined thresholds, ensuring that potential hallucinations can be identified before documents are used for historical analysis. Development of such automated quality control systems is a future research direction.

\subsubsection{Training Data Bias}

The synthetic dataset used for GAN training (derived from relatively clean ANRI documents with multi-stage synthetic degradation pipelines) does not fully replicate natural degradation of VOC-era documents (1650--1800) due to characteristic differences: (1) synthetic degradation (bleed-through, texture noise, blur) does not capture the full complexity of historical degradation such as iron gall ink corrosion with local pH changes, (2) synthetic aging patterns do not fully replicate natural degradation over centuries, and (3) archival storage condition variations (humidity, temperature, light exposure) affecting real document degradation.

To mitigate this bias, future research can explore several strategies: fine-tuning on real degraded ANRI documents for domain adaptation, few-shot learning with minimal examples from target domain, and corrective feedback integration involving paleography experts in the loop for challenging cases.

\subsubsection{Recommendations for Practitioners}

For archival managers, it is recommended to: (1) use restored images as starting points for analysis, not as final truth, (2) always verify critical results by comparing with original documents, (3) document model information (version, configuration, restoration date) in archival records, and (4) implement dual archival protocols to maintain original document integrity. The current limitation is that the system does not yet provide automatic confidence scores, so quality assessment requires visual inspection by experts.

For system developers who will implement this method in production, it is recommended to: (1) implement explainability tools such as attention maps for restoration process transparency, (2) develop uncertainty estimation for each restoration prediction, (3) maintain model versioning systems for reproducibility, (4) conduct periodic bias \textit{audits} to ensure fairness and accuracy, and (5) integrate hallucination detection mechanisms for quality control.

% ============================================================ 
% VII. CONCLUSION
% ============================================================ 
\section{Conclusion}

This research proposes an HTR-oriented GAN-based document restoration framework that addresses limitations of conventional methods in optimizing text readability for automatic recognition systems. The proposed method integrates two main contributions: (1) a \textit{frozen HTR recognizer} strategy providing stable text-aware gradients without joint training instability, proven to be a key success factor of the system through ablation study ($p<0.001$), and (2) multi-component loss function optimization identifying the 4-component configuration (pixel, adversarial, perceptual, CTC) as the optimal balance between visual quality and HTR readability. Exploration of dual-modal discriminator architecture shows statistically insignificant marginal contribution ($\Delta$PSNR +0.28~dB, $p>0.05$), confirming that training protocol and frozen recognizer strategy are more dominant than discriminator architecture complexity.

Quantitative evaluation on 712 test images from a realistic semi-synthetic dataset based on 16th-18th century ANRI paleographic documents demonstrates CER reduction from 83.4\% (degraded condition) to 34.9\% (58.2\% relative reduction, $p<0.001$), approaching theoretical performance on clean images (CER 34.1\%), while maintaining high visual quality (PSNR 30.74~dB, SSIM 0.987). Qualitative validation on 15 authentic ANRI manuscripts confirms effective generalization to real historical degradation that cannot be synthetically replicated. Systematic ablation study validates that the 4-component loss combination achieves optimality in multi-objective space, with the important finding that recognizer feature loss is counterproductive and should be avoided.

The main finding of this research is that explicit integration of HTR objectives through frozen recognizer during GAN training produces outputs optimized for machine readability while maintaining visual quality, with practical implications for large-scale historical archival digitization projects. Limitations include challenges with extreme degradation (contrast $<$30, very heavy fading), complex paleographic ligatures (62.9\% of failure cases), and restriction to Latin paleographic script. Future research will focus on paleographic ligature preservation, historical chemical degradation modeling (iron gall ink corrosion), layered document separation, low-contrast robustness improvement, and transfer-learning validation on multi-language and multi-script corpora to confirm cross-domain generalization.

% ============================================================ 
% REFERENCES
% ============================================================ 
\bibliographystyle{IEEEtran}
\begin{thebibliography}{99}

\bibitem{otsu1979threshold}
N. Otsu, ``A threshold selection method from gray-level histograms,'' \textit{IEEE Trans. Syst., Man, Cybern.}, vol. 9, no. 1, pp. 62--66, Jan. 1979.

\bibitem{sauvola2000adaptive}
J. Sauvola and M. Pietikäinen, ``Adaptive document image binarization,'' \textit{Pattern Recognit.}, vol. 33, no. 2, pp. 225--236, Feb. 2000.

\bibitem{niblack1986introduction}
W. Niblack, \textit{An Introduction to Digital Image Processing}. Englewood Cliffs, NJ, USA: Prentice-Hall, 1986.

\bibitem{wolf2004extraction}
C. Wolf, J.-M. Jolion, and F. Chassaing, ``Text localization, enhancement and binarization in multimedia documents,'' in \textit{Proc. 16th Int. Conf. Pattern Recognit.}, vol. 2, 2002, pp. 1037--1040.

\bibitem{tensmeyer2017document}
C. Tensmeyer and T. Martinez, ``Document image binarization with fully convolutional neural networks,'' in \textit{Proc. 14th IAPR Int. Conf. Document Anal. Recognit.}, vol. 1, 2017, pp. 99--104.

\bibitem{zhao2019document}
J. Zhao, C. Shi, F. Jia, Y. Wang, and B. Xiao, ``Document image binarization with cascaded generators of conditional generative adversarial networks,'' \textit{Pattern Recognit.}, vol. 96, p. 106968, Dec. 2019.

\bibitem{souibgui2020degan}
M. A. Souibgui, S. Khamekhem Jemni, Y. Kessentini, C. Fornés, A. Fornés, and C. Ogier, ``DE-GAN: A conditional generative adversarial network for document enhancement,'' \textit{IEEE Trans. Pattern Anal. Mach. Intell.}, vol. 44, no. 3, pp. 1180--1191, Mar. 2022.

\bibitem{souibgui2021enhance}
M. A. Souibgui, S. Khamekhem Jemni, Y. Kessentini, and A. Fornés, ``Enhance to read better: A multi-task adversarial network for handwritten document image enhancement,'' \textit{Pattern Recognit.}, vol. 123, p. 108368, Mar. 2022.

\bibitem{ronneberger2015unet}
O. Ronneberger, P. Fischer, and T. Brox, ``U-Net: Convolutional networks for biomedical image segmentation,'' in \textit{Proc. Int. Conf. Med. Image Comput. Comput.-Assist. Interv.}, 2015, pp. 234--241.

\bibitem{johnson2016perceptual}
J. Johnson, A. Alahi, and L. Fei-Fei, ``Perceptual losses for real-time style transfer and super-resolution,'' in \textit{Proc. Eur. Conf. Comput. Vis.}, 2016, pp. 694--711.

\bibitem{wang2004ssim}
Z. Wang, A. C. Bovik, H. R. Sheikh, and E. P. Simoncelli, ``Image quality assessment: From error visibility to structural similarity,'' \textit{IEEE Trans. Image Process.}, vol. 13, no. 4, pp. 600--612, Apr. 2004.

\bibitem{marti2002iam}
U.-V. Marti and H. Bunke, ``The IAM-database: An English sentence database for offline handwriting recognition,'' \textit{Int. J. Document Anal. Recognit.}, vol. 5, no. 1, pp. 39--46, Nov. 2002.

\bibitem{perlin1985image}
K. Perlin, ``An image synthesizer,'' \textit{ACM SIGGRAPH Comput. Graph.}, vol. 19, no. 3, pp. 287--296, Jul. 1985.

\bibitem{isola2017image}
P. Isola, J.-Y. Zhu, T. Zhou, and A. A. Efros, ``Image-to-image translation with conditional adversarial networks,'' in \textit{Proc. IEEE Conf. Comput. Vis. Pattern Recognit.}, 2017, pp. 1125--1134.

\bibitem{zhu2017unpaired}
J.-Y. Zhu, T. Park, P. Isola, and A. A. Efros, ``Unpaired image-to-image translation using cycle-consistent adversarial networks,'' in \textit{Proc. IEEE Int. Conf. Comput. Vis.}, 2017, pp. 2223--2232.

\bibitem{souibgui2022docent}
M. A. Souibgui, Y. Kessentini, and A. Fornés, ``DocEnTr: An end-to-end document image enhancement transformer,'' \textit{Pattern Recognit.}, vol. 142, p. 109710, Oct. 2023.

\bibitem{shi2016crnn}
B. Shi, X. Bai, and C. Yao, ``An end-to-end trainable neural network for image-based sequence recognition and its application to scene text recognition,'' \textit{IEEE Trans. Pattern Anal. Mach. Intell.}, vol. 39, no. 11, pp. 2298--2304, Nov. 2017.

\bibitem{bluche2017scan}
T. Bluche, J. Louradour, and R. Messina, ``Scan, attend and read: End-to-end handwritten paragraph recognition with MDLSTM attention,'' in \textit{Proc. 14th IAPR Int. Conf. Document Anal. Recognit.}, vol. 1, 2017, pp. 1050--1055.

\bibitem{michael2019evaluating}
J. Michael, R. Labahn, T. Grüning, and J. Zöllner, ``Evaluating sequence-to-sequence models for handwritten text recognition,'' in \textit{Proc. 15th Int. Conf. Document Anal. Recognit.}, 2019, pp. 1286--1293.

\bibitem{sanchez2016icfhr}
J. A. Sánchez, V. Romero, A. H. Toselli, and E. Vidal, ``ICFHR2016 competition on handwritten text recognition on the READ dataset,'' in \textit{Proc. 15th Int. Conf. Frontiers Handwriting Recognit.}, 2016, pp. 630--635.

\bibitem{sanchez2017icdar}
J. A. Sánchez, V. Romero, A. H. Toselli, and E. Vidal, ``ICDAR2017 competition on handwritten text recognition on the READ dataset,'' in \textit{Proc. 14th IAPR Int. Conf. Document Anal. Recognit.}, vol. 1, 2017, pp. 1383--1388.

\bibitem{diaz2021rethinking}
D. Díaz, A. H. Toselli, E. Vidal, and A. Fornés, ``Rethinking optical character recognition with transformers,'' \textit{Pattern Recognit.}, vol. 122, p. 108301, Feb. 2022.

\bibitem{goodfellow2016improved}
I. Goodfellow, J. Pouget-Abadie, M. Mirza, B. Xu, D. Warde-Farley, S. Ozair, A. Courville, and Y. Bengio, ``Generative adversarial nets,'' in \textit{Proc. Adv. Neural Inf. Process. Syst.}, 2014, pp. 2672--2680.

\bibitem{zhang2018residual}
Y. Zhang, Y. Tian, Y. Kong, B. Zhong, and Y. Fu, ``Residual dense network for image super-resolution,'' in \textit{Proc. IEEE Conf. Comput. Vis. Pattern Recognit.}, 2018, pp. 2472--2481.

\bibitem{woo2018cbam}
S. Woo, J. Park, J.-Y. Lee, and I. S. Kweon, ``CBAM: Convolutional block attention module,'' in \textit{Proc. Eur. Conf. Comput. Vis.}, 2018, pp. 3--19.

\bibitem{lin2017fpn}
T.-Y. Lin, P. Dollár, R. Girshick, K. He, B. Hariharan, and S. Belongie, ``Feature pyramid networks for object detection,'' in \textit{Proc. IEEE Conf. Comput. Vis. Pattern Recognit.}, 2017, pp. 2117--2125.

\bibitem{chen2017deeplab}
L.-C. Chen, G. Papandreou, I. Kokkinos, K. Murphy, and A. L. Yuille, ``DeepLab: Semantic image segmentation with deep convolutional nets, atrous convolution, and fully connected CRFs,'' \textit{IEEE Trans. Pattern Anal. Mach. Intell.}, vol. 40, no. 4, pp. 834--848, Apr. 2018.

\bibitem{oktay2018attention}
O. Oktay et al., ``Attention U-Net: Learning where to look for the pancreas,'' in \textit{Proc. Med. Imag. Deep Learn.}, 2018, pp. 1--10.

\bibitem{kirkpatrick2017overcoming}
J. Kirkpatrick et al., ``Overcoming catastrophic forgetting in neural networks,'' \textit{Proc. Natl. Acad. Sci. USA}, vol. 114, no. 13, pp. 3521--3526, Mar. 2017.

\bibitem{souibgui2022}
M. A. Souibgui et al., ``Enhancing handwritten text recognition with N-gram sequence alignment: A hybrid loss approach for improved accuracy,'' \textit{Pattern Recognit.}, vol. 125, p. 108515, May 2022.

\bibitem{kang2021pay}
L. Kang, P. Riba, M. Rusiñol, A. Fornés, and M. Villegas, ``Pay attention to what you read: Non-recurrent handwritten text-line recognition,'' \textit{Pattern Recognit.}, vol. 129, p. 108766, Sep. 2022.

\bibitem{he2016deep}
K. He, X. Zhang, S. Ren, and J. Sun, ``Deep residual learning for image recognition,'' in \textit{Proc. IEEE Conf. Comput. Vis. Pattern Recognit.}, 2016, pp. 770--778.

\bibitem{hochreiter1997lstm}
S. Hochreiter and J. Schmidhuber, ``Long short-term memory,'' \textit{Neural Comput.}, vol. 9, no. 8, pp. 1735--1780, Nov. 1997.

\bibitem{vaswani2017attention}
A. Vaswani et al., ``Attention is all you need,'' in \textit{Proc. Adv. Neural Inf. Process. Syst.}, 2017, pp. 5998--6008.

\bibitem{pratikakis2016hdibco}
I. Pratikakis, K. Zagoris, G. Barlas, and B. Gatos, ``ICFHR2016 handwritten document image binarization contest (H-DIBCO 2016),'' in \textit{Proc. 15th Int. Conf. Frontiers Handwriting Recognit.}, 2016, pp. 619--623.

\end{thebibliography}

% ============================================================ 
% ACKNOWLEDGMENTS
% ============================================================ 
\section*{Acknowledgments}

The authors would like to thank the National Archives of the Republic of Indonesia (ANRI) for providing access to historical document collections and computational infrastructure that was invaluable for this research. Special appreciation is directed to Bandung Institute of Technology (ITB), particularly the Smart System Study Program at the School of Electrical Engineering and Informatics for providing the research environment that supported the development of this research. We also extend our gratitude to the Ministry of Communication and Informatics of the Republic of Indonesia (Komdigi) for scholarship support that made this research possible.

% ============================================================ 
% BIOGRAFI
% ============================================================ 

\vfill\pagebreak

% \noindent
% \hfill
\begin{minipage}[t]{0.48\textwidth}
\begin{IEEEbiography}[{
\includegraphics[width=1in,height=1.25in,clip,keepaspectratio]{/home/lambda_one/tesis/GAN-HTR-ORI/docRestoration/Paper/penulis/fotoJatniko.png}}]{Jatniko Nur Mutaqin}
is currently pursuing a master's degree at Bandung Institute of Technology (ITB) in the Smart System Study Program, School of Electrical Engineering and Informatics. He works as an AI Engineer at the National Archives of the Republic of Indonesia (ANRI), where his practical experience in historical document preservation and digitization inspires his academic research. His research interests include document image analysis, \textit{generative adversarial networks} (GAN), \textit{optical character recognition} and \textit{handwritten text recognition} systems, and artificial intelligence applications for cultural heritage and historical document preservation.
\end{IEEEbiography}
\end{minipage}

\vspace{1em}

\noindent
\hfill
\begin{minipage}[t]{0.48\textwidth}
\begin{IEEEbiography}[{
\includegraphics[width=1in,height=1.25in,clip,keepaspectratio]{/home/lambda_one/tesis/GAN-HTR-ORI/docRestoration/Paper/penulis/fotoBhaskara.jpg}}]{I Gusti Bagus Baskara Nugraha}
(Ir., S.T., M.T., Ph.D.) is an academic and researcher affiliated with Bandung Institute of Technology (ITB), specifically at the School of Electrical Engineering and Informatics. His expertise lies in Management Information Systems, with research topics including smart city platforms, blockchain-based systems, machine learning, and energy management systems. He has contributed to various academic papers and conferences, with leading research in blockchain, smart city technology, and artificial intelligence applications.
\end{IEEEbiography}
\end{minipage}

\end{document}